\documentclass[12pt,a4paper,oneside]{report}

\usepackage[T1]{fontenc}
\usepackage[utf8]{inputenc}
\usepackage{amsmath}
\let\Bbbk\relax
\usepackage{amssymb}
\usepackage{newtxtext,newtxmath}
\usepackage[a4paper,top=1.1in,bottom=1.1in,left=1.25in,right=1.1in]{geometry}
\usepackage{graphicx}
\usepackage{booktabs}
\usepackage{longtable}
\usepackage{array}
\usepackage{siunitx}
\usepackage{xcolor}
\usepackage{titlesec}
\usepackage[font=small,labelfont=bf,width=0.92\textwidth]{caption}
\usepackage{enumitem}
\usepackage{tocloft}
\usepackage{microtype}
\usepackage{emptypage}
\usepackage{tikz}
\usetikzlibrary{arrows.meta,positioning,calc,shapes.geometric,fit,backgrounds}
\usepackage[american,siunitx]{circuitikz}
\usepackage[hidelinks,bookmarksnumbered]{hyperref}

\graphicspath{{../../results/figures/}}

\definecolor{accent}{RGB}{0,86,140}
\definecolor{rule}{RGB}{120,120,120}

\titleformat{\chapter}[display]
  {\normalfont\huge\bfseries\raggedright}
  {\normalfont\normalsize\color{accent}\MakeUppercase{\chaptertitlename}\ \thechapter}
  {10pt}
  {\color{black}}
\titleformat{name=\chapter,numberless}[display]
  {\normalfont\huge\bfseries\raggedright}{}{0pt}{\color{black}}
\titlespacing*{\chapter}{0pt}{0pt}{28pt}

\titleformat{\section}{\normalfont\large\bfseries}{\thesection}{0.7em}{}
\titleformat{\subsection}{\normalfont\normalsize\bfseries}{\thesubsection}{0.7em}{}
\titleformat{\subsubsection}{\normalfont\normalsize\bfseries\itshape}{\thesubsubsection}{0.7em}{}
\titlespacing*{\section}{0pt}{16pt}{6pt}
\titlespacing*{\subsection}{0pt}{12pt}{4pt}

\pagestyle{empty}
\makeatletter\let\ps@plain\ps@empty\makeatother
\linespread{1.06}\selectfont

\setlength{\parskip}{0.45em}
\setlength{\parindent}{0pt}
\renewcommand{\arraystretch}{1.15}
\sisetup{per-mode=symbol, detect-weight=true}
\DeclareSIUnit\ohm{\ensuremath{\Omega}}
\DeclareSIUnit\celsius{\ensuremath{^{\circ}}C}
\DeclareSIUnit\degree{\ensuremath{^{\circ}}}
\DeclareSIUnit\percent{\ensuremath{\%}}
\DeclareSIUnit\electronvolt{eV}
\DeclareSIPrefix\micro{\ensuremath{\mu}}{-6}

\renewcommand{\cftchapfont}{\bfseries}
\renewcommand{\cftchappagefont}{\bfseries}
\setlength{\cftbeforechapskip}{6pt}

\newcommand{\term}[1]{\textbf{#1}}
\newcommand{\vpv}{v_{\mathrm{pv}}}
\newcommand{\ipv}{i_{\mathrm{pv}}}
\newcommand{\Vmp}{V_{\mathrm{mp}}}
\newcommand{\Imp}{I_{\mathrm{mp}}}
\newcommand{\Pmp}{P_{\mathrm{mp}}}
\newcommand{\Voc}{V_{\mathrm{oc}}}
\newcommand{\Isc}{I_{\mathrm{sc}}}

\begin{document}

\begin{titlepage}
\centering

\IfFileExists{logo.png}{%
  \vspace*{0.35in}
  \includegraphics[width=0.52\textwidth]{logo.png}\par
  \vspace{0.45in}
}{%
  \vspace*{1.2in}
}

\vspace{0.15in}

{\Huge\bfseries Solar Panel Maximum Power\\[0.25em] Point Tracking as a\\[0.25em] Control Problem\par}

\vspace{0.4in}

{\large\itshape Feedback control for operation near maximum power\\ under changing irradiance\par}

\vspace{0.5in}

\vspace{0.75in}

{\large\bfseries Roshan\par}

\vspace{0.3in}

{\large A project report on the modelling, control design, simulation\\ and physical realisation of a photovoltaic maximum power\\ point tracking system\par}

\vfill

{\large \today\par}
\end{titlepage}

\pagenumbering{roman}
\setcounter{page}{1}
\pagestyle{empty}

\tableofcontents
\clearpage
\listoffigures
\clearpage
\listoftables
\clearpage

\chapter*{Nomenclature and Abbreviations}
\addcontentsline{toc}{chapter}{Nomenclature and Abbreviations}
\markboth{Nomenclature and Abbreviations}{}

Every abbreviation used anywhere in this report is expanded here, and every
symbol that appears in an equation is defined here as well as at its first use.

\section*{Abbreviations}

\begin{longtable}{@{}p{0.16\textwidth}p{0.80\textwidth}@{}}
\toprule
\textbf{Abbreviation} & \textbf{Full form and meaning} \\
\midrule
\endfirsthead
\toprule
\textbf{Abbreviation} & \textbf{Full form and meaning} \\
\midrule
\endhead
\bottomrule
\endfoot

AC & Alternating Current. An electric current whose direction reverses periodically. \\
ADC & Analogue-to-Digital Converter. A circuit that converts a continuous voltage into a number a processor can read. \\
CCM & Continuous Conduction Mode. Operating condition of a switching converter in which the inductor current never falls to zero within a switching period. \\
CVD & Colour Vision Deficiency. Reduced ability to distinguish certain colours; considered when choosing figure colours. \\
DC & Direct Current. An electric current of constant direction. \\
DCM & Discontinuous Conduction Mode. The complementary case to CCM, in which inductor current reaches zero and stays there for part of the switching period. \\
DSP & Digital Signal Processor. A microprocessor specialised for real-time arithmetic, used here as the control processor. \\
EMI & Electromagnetic Interference. Unwanted radiated or conducted electrical noise. \\
ESC & Extremum Seeking Control. A feedback method that drives a system to the extremum of a measured performance function without knowing that function. \\
ESR & Equivalent Series Resistance. The effective resistance in series with an ideal capacitor or inductor. \\
GaN & Gallium Nitride. A wide-bandgap semiconductor used for fast, low-loss power transistors. \\
IGBT & Insulated Gate Bipolar Transistor. A power switching device suited to higher voltages. \\
INC & Incremental Conductance. An MPPT algorithm that compares the incremental conductance $\mathrm{d}I/\mathrm{d}V$ with the instantaneous conductance $-I/V$. \\
I--V & Current--Voltage. Refers to the current-versus-voltage characteristic curve of a solar device. \\
LC & Inductor--Capacitor. Refers to the resonant circuit formed by an inductor and a capacitor. \\
LHP & Left Half Plane. The region of the complex plane with negative real part; poles there correspond to stable modes. \\
LPV & Linear Parameter Varying. A model that is linear for each fixed value of a parameter, but whose coefficients depend on that parameter. \\
LTI & Linear Time Invariant. A system whose behaviour is linear and does not change with time. \\
MOSFET & Metal Oxide Semiconductor Field Effect Transistor. The transistor used as the controlled switch in the converter. \\
MPP & Maximum Power Point. The single operating point on a photovoltaic I--V curve at which output power is greatest. \\
MPPT & Maximum Power Point Tracking. The control function that keeps a photovoltaic source operating at or near its MPP. \\
NOCT & Nominal Operating Cell Temperature. The cell temperature reached under a defined reference environment, used to estimate cell temperature in the field. \\
ODE & Ordinary Differential Equation. \\
P\&O & Perturb and Observe. An MPPT algorithm that repeatedly perturbs the operating point and keeps moving in whichever direction increased power. \\
PCB & Printed Circuit Board. \\
PI & Proportional--Integral. A controller whose output is a weighted sum of the error and the integral of the error. \\
PID & Proportional--Integral--Derivative. A PI controller with an added derivative term; not used here. \\
PM & Phase Margin. The additional phase lag that would make a stable feedback loop marginally stable. \\
PS & Physical Signal. In Simscape, a directed signal carrying a physical quantity with units, as distinct from an electrical connection. \\
P--V & Power--Voltage. Refers to the power-versus-voltage characteristic curve of a solar device. \\
PV & Photovoltaic. Relating to the direct conversion of light into electricity. \\
PWM & Pulse Width Modulation. Encoding a continuous command as the duty ratio of a fixed-frequency rectangular wave. \\
RCC & Ripple Correlation Control. An MPPT method that correlates the converter's own switching ripple in power and voltage to estimate the power gradient. \\
RHP & Right Half Plane. The region of the complex plane with positive real part; a zero there fundamentally limits achievable bandwidth. \\
RK4 & Fourth-order Runge--Kutta. A numerical method for integrating ordinary differential equations. \\
RMS & Root Mean Square. \\
SiC & Silicon Carbide. A wide-bandgap semiconductor used for efficient power devices. \\
STC & Standard Test Conditions. \SI{1000}{\watt\per\square\metre} irradiance, \SI{25}{\celsius} cell temperature, air mass 1.5 spectrum. \\
\bottomrule
\end{longtable}

\section*{Symbols}

\begin{longtable}{@{}p{0.16\textwidth}p{0.62\textwidth}p{0.14\textwidth}@{}}
\toprule
\textbf{Symbol} & \textbf{Meaning} & \textbf{Unit} \\
\midrule
\endfirsthead
\toprule
\textbf{Symbol} & \textbf{Meaning} & \textbf{Unit} \\
\midrule
\endhead
\bottomrule
\endfoot

$a$ & Modified ideality factor, $a = n N_s k T_c / q$ & \si{\volt} \\
$C_{\mathrm{pv}}$ & Input capacitance across the photovoltaic terminals & \si{\farad} \\
$C_o$ & Output capacitance at the DC bus & \si{\farad} \\
$d$ & Duty ratio, the fraction of each switching period the main switch conducts & -- \\
$D$ & Steady-state value of the duty ratio & -- \\
$E_g$ & Semiconductor bandgap energy & \si{\electronvolt} \\
$f_{sw}$ & Switching frequency & \si{\hertz} \\
$f_{c}$ & Loop gain crossover frequency & \si{\hertz} \\
$G$ & Plane-of-array irradiance & \si{\watt\per\square\metre} \\
$g$ & Slope of the I--V curve, $g = \mathrm{d}I/\mathrm{d}V$ & \si{\siemens} \\
$i_L$ & Inductor current & \si{\ampere} \\
$\ipv$ & Photovoltaic array terminal current & \si{\ampere} \\
$I_0$ & Diode reverse saturation current & \si{\ampere} \\
$I_{ph}$ & Photocurrent, the light-generated current & \si{\ampere} \\
$\Isc$ & Short-circuit current & \si{\ampere} \\
$\Imp$ & Current at the maximum power point & \si{\ampere} \\
$k$ & Boltzmann constant, $\num{1.380649e-23}$ & \si{\joule\per\kelvin} \\
$K_p$, $K_i$ & Proportional and integral gains of a PI controller & -- \\
$L$ & Boost inductance & \si{\henry} \\
$n$ & Diode ideality factor & -- \\
$N_s$ & Number of cells in series within one module & -- \\
$N_{sm}$, $N_{ps}$ & Modules in series, strings in parallel & -- \\
$P$ & Photovoltaic output power & \si{\watt} \\
$\Pmp$ & Power at the maximum power point & \si{\watt} \\
$q$ & Elementary charge, $\num{1.602176634e-19}$ & \si{\coulomb} \\
$r_d$ & Dynamic resistance of the array, $r_d = -\mathrm{d}V/\mathrm{d}I$ & \si{\ohm} \\
$R_L$ & Equivalent series resistance of the boost inductor & \si{\ohm} \\
$R_s$ & Lumped series resistance of the photovoltaic model & \si{\ohm} \\
$R_{sh}$ & Shunt (parallel) resistance of the photovoltaic model & \si{\ohm} \\
$T_c$ & Cell temperature & \si{\kelvin} \\
$T_s$ & Sampling period of a discrete controller & \si{\second} \\
$V_{dc}$ & DC bus voltage & \si{\volt} \\
$\vpv$ & Photovoltaic array terminal voltage & \si{\volt} \\
$\Voc$ & Open-circuit voltage & \si{\volt} \\
$\Vmp$ & Voltage at the maximum power point & \si{\volt} \\
$\alpha_{Isc}$ & Temperature coefficient of short-circuit current & \si{\ampere\per\kelvin} \\
$\beta_{Voc}$ & Temperature coefficient of open-circuit voltage & \si{\volt\per\kelvin} \\
$\zeta$ & Damping ratio of a second-order mode & -- \\
$\eta_{\mathrm{MPPT}}$ & MPPT efficiency, harvested energy divided by available energy & -- \\
$\omega_n$ & Undamped natural frequency & \si{\radian\per\second} \\
\bottomrule
\end{longtable}


\clearpage
\pagenumbering{arabic}
\setcounter{page}{1}

\chapter{Project Statement and Scope}
\label{ch:statement}

\section{The project as originally stated}

The project was defined as follows.

\begin{quote}
\textbf{Project Title:} Solar Panel Maximum Power Point Tracking as a Control
Problem.

Core idea is to use feedback control to maintain operations near maximum power
despite changing irradiance. In MATLAB, need to model PV model, MPPT model, and
control analysis. In physical model, need solar panel and DC-DC model.
\end{quote}

Everything in this report follows from that statement. This first chapter
unpacks it: what each phrase means, what a complete answer has to contain, and
how the work was organised to deliver it.

\section{Reading the statement}

\subsection{``Maximum power point''}

A solar panel is neither a voltage source nor a current source. Its terminal
current depends nonlinearly on its terminal voltage, and the product of the two,
the delivered power, has a single interior maximum. That maximum is the
\term{maximum power point} (MPP). Operating anywhere else wastes energy that the
sunlight already delivered to the panel.

\subsection{``Despite changing irradiance''}

The location of the MPP is not fixed. It moves when the \term{irradiance}, the
optical power per unit area reaching the panel, changes because of clouds, time
of day or shading. It also moves when the cell temperature changes. So the
problem is not to find a number once; it is to \emph{track} a moving target
continuously, in real time, using only measurements of terminal voltage and
terminal current.

\subsection{``As a control problem''}

This is the phrase that sets the character of the whole project. Much of the
published work on MPPT treats it as an \emph{algorithm} question: which search
rule finds the peak fastest. That framing hides the fact that the search rule
sits on top of a physical plant, a switching power converter, whose dynamics
decide whether the search rule can work at all.

Treating MPPT as a control problem means asking the questions a control engineer
would ask:

\begin{itemize}[leftmargin=1.4em]
  \item What exactly is the plant, and what is its transfer function?
  \item How do the plant dynamics change as the operating point moves along the
        I--V curve and as irradiance changes?
  \item What is the control input, what is the measured output, and what is the
        error signal?
  \item What stability margins does the closed loop have, and are they
        maintained over the whole operating envelope?
  \item What fundamental limits, right half plane zeros, resonances, sampling
        delay, bound the achievable bandwidth?
  \item Given those limits, how fast can a tracker legitimately be, and what
        does it cost in steady-state accuracy?
\end{itemize}

A central result of this project, developed in
Chapter~\ref{ch:converter} and Chapter~\ref{ch:control}, is that these questions
have sharp and somewhat surprising answers, and that the answers explain the
behaviour of the classical MPPT algorithms rather than merely sitting alongside
them.

\subsection{``In MATLAB, need to model PV model, MPPT model, and control analysis''}

Three deliverables on the analysis side:

\begin{enumerate}[leftmargin=1.6em]
  \item A \term{photovoltaic model}: a mathematical description of the panel
        accurate enough that its curvature, not just its three datasheet points,
        can be trusted. Chapter~\ref{ch:pvmodel}.
  \item An \term{MPPT model}: implementations of the tracking algorithms as
        proper discrete-time controllers with defined sample rates and states,
        so that they can be compared fairly. Chapter~\ref{ch:mppt}.
  \item A \term{control analysis}: linearisation, frequency response, stability
        margins and robustness across the operating envelope.
        Chapters~\ref{ch:converter} and~\ref{ch:control}.
\end{enumerate}

\subsection{``In physical model, need solar panel and DC-DC model''}

A separate deliverable: a model of the actual hardware, not an idealised
abstraction. The panel appears as a physical two-terminal electrical device, and
the \term{DC--DC converter}, the switching circuit that transforms one direct
voltage into another, appears as a real circuit with a transistor that switches
at \SI{20}{\kilo\hertz}, a diode, an inductor with resistance and capacitors.
Chapter~\ref{ch:physical} covers this, together with the hardware required to
build it.

\section{What a complete answer requires}

Putting those readings together, a complete project needs all of the following,
and each is delivered in the chapter named.

\begin{longtable}{@{}L{0.40\textwidth}L{0.42\textwidth}L{0.12\textwidth}@{}}
\toprule
\textbf{Requirement} & \textbf{How it is met} & \textbf{Chapter} \\
\midrule
\endfirsthead
\toprule
\textbf{Requirement} & \textbf{How it is met} & \textbf{Chapter} \\
\midrule
\endhead
\bottomrule
\endfoot
A panel model valid away from the datasheet points & Five-parameter single-diode model with De Soto translation, parameters identified from five simultaneous datasheet constraints & \ref{ch:pvmodel} \\
Ground truth to score trackers against & Exact maximum power point by root solve on the analytic power gradient & \ref{ch:pvmodel} \\
A converter model suitable for control design & State-space averaged boost converter, linearised analytically, Jacobian verified against finite differences & \ref{ch:converter} \\
Understanding of how the plant changes & Operating envelope sweep in irradiance and along the I--V curve & \ref{ch:converter} \\
A controller with quantified margins & Cascaded current and voltage loops, margins computed over the envelope, compared against a single-loop baseline & \ref{ch:control} \\
MPPT algorithms implemented fairly & Common tracker interface, identical sample rate and identical inner control loops & \ref{ch:mppt} \\
An objective comparison metric & Dynamic MPPT efficiency on EN 50530-style irradiance ramps & \ref{ch:benchmark} \\
Reproducibility and correctness evidence & Scripted workflow and 27 automated unit tests & \ref{ch:software} \\
A physical circuit model & Switched Simscape model with a custom photovoltaic component, cross-validated against the averaged model & \ref{ch:physical} \\
A route to hardware & Component selection, ratings, sensing, protection, and a staged build plan & \ref{ch:physical} \\
\end{longtable}

\section{System chosen for study}

All results in this report are for one concrete system, so that numbers can be
compared directly from chapter to chapter.

\begin{table}[htbp]
\centering
\caption{The system under study.}
\label{tab:system}
\begin{tabular}{@{}lll@{}}
\toprule
\textbf{Item} & \textbf{Value} & \textbf{Note} \\
\midrule
Module & Kyocera KC200GT & \SI{200}{\watt}, 54 cells, multicrystalline silicon \\
Array & 6 series $\times$ 2 parallel & \SI{2402}{\watt} at STC \\
Array MPP at STC & \SI{157.8}{\volt}, \SI{15.22}{\ampere} & \\
Converter & Boost & PV side to DC bus \\
Inductance $L$ & \SI{1.5}{\milli\henry} & \SI{20}{\percent} current ripple \\
Input capacitance $C_{\mathrm{pv}}$ & \SI{100}{\micro\farad} & sets PV voltage ripple \\
Inductor resistance $R_L$ & \SI{0.15}{\ohm} & \\
Switching frequency $f_{sw}$ & \SI{20}{\kilo\hertz} & \\
DC bus $V_{dc}$ & \SI{400}{\volt} & treated as stiff \\
\bottomrule
\end{tabular}
\end{table}

The KC200GT was chosen because its datasheet is widely used as a modelling
benchmark in the literature~\cite{villalva}, which makes the extracted
parameters checkable against published values. The array size, roughly
\SI{2.4}{\kilo\watt} into a \SI{400}{\volt} bus, is representative of the
photovoltaic input stage of a small single-phase string inverter.

\section{The control architecture adopted}

The architecture that emerges from the analysis, and which the rest of the
report develops and justifies, is a three-layer hierarchy separated by
timescale.

\begin{table}[htbp]
\centering
\caption{The three-layer control hierarchy and its timescale separation.}
\label{tab:hierarchy}
\begin{tabular}{@{}llll@{}}
\toprule
\textbf{Layer} & \textbf{Controlled variable} & \textbf{Command it issues} & \textbf{Bandwidth} \\
\midrule
MPPT & PV power (indirectly) & voltage reference $\Vmp^\ast$ & \SI{50}{\hertz} \\
Outer loop & PV voltage $\vpv$ & inductor current reference & \SI{250}{\hertz} \\
Inner loop & inductor current $i_L$ & duty ratio $d$ & \SI{2}{\kilo\hertz} \\
\bottomrule
\end{tabular}
\end{table}

Each layer is roughly eight times slower than the one beneath it. That
separation is what allows each to be designed as though the layer below were
instantaneous and the layer above were frozen, and Chapter~\ref{ch:control}
verifies that the assumption holds by computing margins on the full-order model
rather than on the simplified one used for design.

\chapter{Theoretical Background}
\label{ch:theory}

This chapter defines every concept the later chapters rely on. It is written so
that a reader who knows basic circuit theory and basic calculus can follow the
whole report without consulting anything else.

\section{Photovoltaic conversion}

\subsection{The photovoltaic effect}

A \term{solar cell} is a large-area semiconductor diode. When a photon with
energy greater than the semiconductor \term{bandgap} $E_g$ is absorbed, it lifts
an electron into the conduction band and leaves a hole behind. The built-in
electric field of the p--n junction sweeps these carriers in opposite
directions, producing a current. That current, called the \term{photocurrent}
$I_{ph}$, is very nearly proportional to the incident irradiance $G$.

\subsection{Irradiance and standard test conditions}

\term{Irradiance} is optical power per unit area, measured in
\si{\watt\per\square\metre}. Bright midday sun on a surface facing it is close
to \SI{1000}{\watt\per\square\metre}. Because output depends on light,
temperature and spectrum, modules are rated at \term{Standard Test Conditions}
(STC): irradiance \SI{1000}{\watt\per\square\metre}, cell temperature
\SI{25}{\celsius}, and an air mass 1.5 spectrum.

\subsection{The four datasheet points}

Manufacturers publish, at STC:

\begin{itemize}[leftmargin=1.4em]
  \item $\Isc$, the \term{short-circuit current}: current when the terminals are
        shorted, so the terminal voltage is zero.
  \item $\Voc$, the \term{open-circuit voltage}: voltage when no current is
        drawn.
  \item $\Vmp$ and $\Imp$, the voltage and current at the maximum power point,
        whose product is the rated power $\Pmp$.
  \item $\alpha_{Isc}$ and $\beta_{Voc}$, the \term{temperature coefficients}
        that say how $\Isc$ and $\Voc$ change per kelvin.
\end{itemize}

The \term{fill factor} $FF = \Pmp / (\Voc \Isc)$ measures how square the I--V
curve is; a good crystalline silicon module is around \num{0.75}.

\section{The single-diode model}

\subsection{Circuit and equation}

The standard lumped circuit model of a photovoltaic cell or module places a
current source $I_{ph}$ in parallel with a diode and a shunt resistance
$R_{sh}$, with a series resistance $R_s$ to the terminals. Applying Kirchhoff's
current law gives the \term{single-diode equation}:

\begin{equation}
I = I_{ph} - I_0\!\left[\exp\!\left(\frac{V + I R_s}{a}\right) - 1\right] - \frac{V + I R_s}{R_{sh}} ,
\label{eq:singlediode}
\end{equation}

where

\begin{equation}
a = \frac{n N_s k T_c}{q}
\label{eq:idealityvolts}
\end{equation}

is the \term{modified ideality factor}, expressed in volts. Here $n$ is the
dimensionless diode \term{ideality factor} (unity for an ideal diode, larger
when recombination is significant), $N_s$ is the number of cells in series in
the module, $k$ is the Boltzmann constant, $T_c$ the absolute cell temperature
and $q$ the elementary charge. $I_0$ is the \term{reverse saturation current} of
the diode.

The physical meaning of each element is worth stating plainly:

\begin{itemize}[leftmargin=1.4em]
  \item $I_{ph}$ is the current the light would deliver if nothing internal
        stole it.
  \item The diode term is the recombination current that flows back through the
        junction; it is negligible at low voltage and dominates near $\Voc$.
  \item $R_{sh}$ represents leakage paths across the junction; a small $R_{sh}$
        tilts the flat part of the I--V curve downward.
  \item $R_s$ represents contact, grid and bulk resistance; it rounds the knee
        of the curve.
\end{itemize}

\subsection{Why the equation is implicit}

Equation~\eqref{eq:singlediode} has $I$ on both sides, inside the exponential as
well as outside. There is no elementary closed form for $I(V)$. It can be
written using the Lambert $W$ function, but the argument overflows in double
precision for realistic parameters. This report solves it numerically by a
damped Newton iteration, described in Section~\ref{sec:solver}.

\subsection{Translation to other conditions}

The five reference parameters are identified at STC and then translated to the
actual operating irradiance and temperature. This project uses the De Soto
formulation~\cite{desoto}, which is the basis of the widely used NREL System
Advisor Model:

\begin{align}
a(T_c) &= a_{ref}\,\frac{T_c}{T_{ref}} \label{eq:tr_a}\\
I_{ph}(G,T_c) &= \frac{G}{G_{ref}}\left[I_{ph,ref} + \alpha_{Isc}\,(T_c - T_{ref})\right] \label{eq:tr_iph}\\
E_g(T_c) &= E_{g,ref}\left[1 - \num{2.677e-4}\,(T_c - T_{ref})\right] \label{eq:tr_eg}\\
I_0(T_c) &= I_{0,ref}\left(\frac{T_c}{T_{ref}}\right)^{\!3}\exp\!\left[\frac{E_{g,ref}}{k T_{ref}} - \frac{E_g(T_c)}{k T_c}\right] \label{eq:tr_i0}\\
R_{sh}(G) &= R_{sh,ref}\,\frac{G_{ref}}{G}, \qquad R_s = R_{s,ref} \label{eq:tr_r}
\end{align}

Physically: the thermal voltage rises with temperature; the photocurrent scales
with light and drifts slightly with temperature; the saturation current rises
steeply with temperature because the bandgap narrows; and shunt leakage scales
inversely with illumination.

\subsection{Series and parallel arrays}

For $N_{sm}$ modules in series and $N_{ps}$ strings in parallel, voltages add
along a string and currents add across strings, so the array-level parameters
are

\begin{equation}
\begin{gathered}
a_{arr} = a N_{sm}, \qquad
I_{ph,arr} = I_{ph} N_{ps}, \qquad
I_{0,arr} = I_0 N_{ps},\\[4pt]
R_{s,arr} = R_s \frac{N_{sm}}{N_{ps}}, \qquad
R_{sh,arr} = R_{sh} \frac{N_{sm}}{N_{ps}} .
\end{gathered}
\label{eq:arrayscaling}
\end{equation}

This assumes all modules see the same irradiance and temperature. Partial
shading breaks that assumption and creates multiple local maxima; it is outside
the scope of this project and is noted as future work in
Chapter~\ref{ch:conclusions}.

\section{Quantities that govern the control problem}

\subsection{The power gradient}

Differentiating the output power $P = V I$ with respect to terminal voltage,

\begin{equation}
\frac{\mathrm{d}P}{\mathrm{d}V} = I + V\,\frac{\mathrm{d}I}{\mathrm{d}V} .
\label{eq:dpdv}
\end{equation}

The maximum power point is where this vanishes. Every MPPT algorithm is, in one
way or another, an attempt to drive $\mathrm{d}P/\mathrm{d}V$ to zero using only
terminal measurements.

\subsection{The slope of the I--V curve}

Implicit differentiation of~\eqref{eq:singlediode} gives an exact expression for
the slope, which is used throughout this report rather than a finite difference:

\begin{equation}
\frac{\mathrm{d}I}{\mathrm{d}V}
= -\,\frac{g_d + 1/R_{sh}}{1 + R_s g_d + R_s/R_{sh}},
\qquad
g_d = \frac{I_0}{a}\exp\!\left(\frac{V + I R_s}{a}\right).
\label{eq:didv}
\end{equation}

Here $g_d$ is the \term{incremental conductance of the diode}.

\subsection{Dynamic resistance}

The \term{dynamic resistance} of the array is

\begin{equation}
r_d = -\frac{\mathrm{d}V}{\mathrm{d}I} = -\left(\frac{\mathrm{d}I}{\mathrm{d}V}\right)^{-1}.
\label{eq:rd}
\end{equation}

This single number is the most important quantity in the whole control analysis.
It describes how the array behaves to small perturbations:

\begin{itemize}[leftmargin=1.4em]
  \item Well below $\Vmp$ the curve is nearly flat, $\mathrm{d}I/\mathrm{d}V$ is
        almost zero, and $r_d$ is very large. The array behaves like a
        \emph{current source}.
  \item Well above $\Vmp$ the curve falls steeply, and $r_d$ is small. The array
        behaves like a \emph{voltage source}.
  \item Exactly at the MPP, combining~\eqref{eq:dpdv} and~\eqref{eq:rd} gives
        the clean result
        \begin{equation}
          r_d\big|_{\mathrm{MPP}} = \frac{\Vmp}{\Imp},
          \label{eq:rdmpp}
        \end{equation}
        that is, the dynamic resistance equals the static resistance at that one
        point.
\end{itemize}

Chapter~\ref{ch:converter} shows that $r_d$ appears directly in the damping term
of the converter's characteristic polynomial, so this order-of-magnitude swing
in $r_d$ becomes an order-of-magnitude swing in closed-loop behaviour.

\section{DC--DC conversion}

\subsection{Why a converter is needed at all}

The array delivers maximum power at one particular voltage, which moves with
weather. The load, here a DC bus feeding an inverter, requires a different and
essentially fixed voltage. A \term{DC--DC converter} sits between them and, by
adjusting its \term{duty ratio}, chooses what voltage the array is held at while
still delivering into the fixed bus. The converter is therefore the actuator of
the MPPT control system.

\subsection{The boost converter}

A \term{boost converter} produces an output voltage higher than its input. It
consists of an inductor, a controlled switch to ground, a diode and an output
capacitor. Over one switching period $T_{sw} = 1/f_{sw}$ the switch is closed
for a fraction $d$ of the time, the \term{duty ratio}. In the ideal lossless
steady state,

\begin{equation}
V_{dc} = \frac{\vpv}{1 - D} \qquad\Longleftrightarrow\qquad D = 1 - \frac{\vpv}{V_{dc}} .
\label{eq:boostratio}
\end{equation}

Increasing $d$ therefore \emph{lowers} the PV voltage for a fixed bus voltage,
which is an important sign to keep straight in the control design.

A boost topology is the right choice here because the array MPP voltage,
\SI{157.8}{\volt}, is below the \SI{400}{\volt} bus at every irradiance of
interest, so the converter never needs to buck.

\subsection{Continuous conduction mode}

If the inductor current never reaches zero within a switching period, the
converter is in \term{continuous conduction mode} (CCM) and~\eqref{eq:boostratio}
holds. The peak-to-peak inductor current ripple in CCM is

\begin{equation}
\Delta i_L = \frac{\vpv\, D}{f_{sw} L},
\label{eq:iripple}
\end{equation}

and CCM requires $i_L > \Delta i_L / 2$. Below that the converter enters
\term{discontinuous conduction mode} (DCM), where the small-signal model changes
character. Section~\ref{sec:sizing} checks CCM across the operating range.

\subsection{State-space averaging}

The converter is a switched, and therefore time-varying, system. \term{State
space averaging}~\cite{middlebrook} replaces the two switch states with their
duty-weighted average over one switching period. The result is a continuous,
nonlinear model whose states are the slowly varying averages of the real
waveforms, valid for frequencies well below $f_{sw}$. That model can then be
linearised about an operating point in the usual way. Chapter~\ref{ch:converter}
does this, and Chapter~\ref{ch:physical} checks the result against a genuine
switched circuit simulation.

\section{Feedback control concepts}

\subsection{Plant, controller and loop gain}

The \term{plant} is the physical system to be controlled, here the converter and
array. The \term{controller} computes the input to the plant from the measured
error. Breaking the feedback path and multiplying the transfer functions around
it gives the \term{loop gain} $L(s)$. Closed-loop stability is decided by
$1 + L(s) = 0$.

\subsection{Stability margins}

The \term{gain crossover frequency} $\omega_c$ is where $|L(j\omega_c)| = 1$; it
is the practical bandwidth of the loop. The \term{phase margin} is

\begin{equation}
\mathrm{PM} = 180^\circ + \angle L(j\omega_c),
\end{equation}

the extra phase lag that would drive the loop to the edge of instability. A
design target of at least \SI{45}{\degree} is conventional. \term{Gain margin}
is the corresponding factor by which the gain could rise before instability.

An important subtlety used in Chapter~\ref{ch:control}: if $|L|$ crosses unity
more than once, a single phase margin number is misleading, and the smallest of
the margins at all crossings must be reported. Resonant plants make this
situation common.

\subsection{The proportional--integral controller}

The \term{PI controller} has transfer function

\begin{equation}
C(s) = K_p + \frac{K_i}{s} .
\label{eq:pi}
\end{equation}

The integral term guarantees zero steady-state error to a constant reference,
because any persistent error would integrate without bound. The proportional
term sets the bandwidth. The zero of the controller sits at $\omega_z = K_i/K_p$
and is placed below crossover so that its phase lag has decayed by the time it
matters.

Because the physical actuator saturates, $d$ must lie in $[0,1]$, a bare
integrator can wind up: the integral state grows while the output is stuck at a
limit, causing a long overshoot on recovery. The implementation used here
includes \term{anti-windup} by back-calculation, feeding the difference between
the requested and the delivered output back into the integrator.

\subsection{Cascade control and timescale separation}

\term{Cascade control} nests a fast inner loop inside a slower outer loop. The
inner loop regulates a variable that responds quickly, here inductor current;
the outer loop commands the inner loop's reference. If the inner loop is
substantially faster, the outer loop can be designed as if the inner loop were
an ideal, instantaneous follower. The benefits are that the inner loop linearises
and damps the plant seen by the outer loop, and that current limiting becomes
trivial because the inner reference is a current.

\subsection{Extremum seeking}

Standard feedback drives an error to zero when the desired value is known. In
MPPT the desired voltage is \emph{not} known. \term{Extremum seeking control}
solves this by adding a small sinusoidal \term{dither} to the command, measuring
how the performance output responds at the dither frequency, and using that
correlation as an estimate of the gradient. Formally, if the plant output is
$P(v)$ and the command is $v = \hat v + A\sin\omega t$, then to first order

\begin{equation}
P \approx P(\hat v) + A \sin(\omega t)\, \frac{\mathrm{d}P}{\mathrm{d}v},
\end{equation}

so multiplying the high-pass filtered $P$ by $\sin\omega t$ and low-pass
filtering yields $\tfrac{A}{2}\,\mathrm{d}P/\mathrm{d}v$. Integrating that
estimate performs gradient ascent~\cite{krstic,leyva}. The price is the
permanent dither, which costs a little power in steady state.

\section{Measuring how well a tracker performs}

\subsection{MPPT efficiency}

The agreed figure of merit is \term{MPPT efficiency}, the fraction of available
energy actually harvested over a test window:

\begin{equation}
\eta_{\mathrm{MPPT}} = \frac{\displaystyle\int_{t_0}^{t_1} \vpv(t)\, \ipv(t)\,\mathrm{d}t}
                            {\displaystyle\int_{t_0}^{t_1} \Pmp\big(G(t), T_c(t)\big)\,\mathrm{d}t} .
\label{eq:eta}
\end{equation}

The denominator requires knowing the true maximum at every instant, which is
available in simulation because the model can be solved exactly. This is why
Chapter~\ref{ch:pvmodel} insists on an exact MPP solver rather than an
approximation.

\subsection{Static and dynamic efficiency}

\term{Static MPPT efficiency} is measured at constant irradiance and penalises
only steady-state error and ripple. \term{Dynamic MPPT efficiency} is measured
while irradiance changes and additionally penalises lag. The European standard
EN 50530~\cite{en50530} defines trapezoidal irradiance ramps for exactly this
purpose; Chapter~\ref{ch:benchmark} uses profiles of that form.

\chapter{Modelling the Photovoltaic Array}
\label{ch:pvmodel}

\section{What the model has to be good at}

For an MPPT study the model must do more than reproduce the three datasheet
points. The controller lives on the \emph{curvature} of the characteristic: the
plant gain is set by the dynamic resistance $r_d$, which is a derivative, and
the MPPT layer operates on $\mathrm{d}P/\mathrm{d}V$, which is another
derivative. A model that passes through $\Isc$, $\Voc$ and the MPP but has the
wrong shape between them would give plausible-looking power numbers and
completely wrong control conclusions.

Three requirements follow:

\begin{enumerate}[leftmargin=1.6em]
  \item The five parameters must be identified from constraints that pin the
        \emph{slope} as well as the position of the curve.
  \item Derivatives must be computed analytically, not by finite differences of
        a numerically solved current.
  \item The maximum power point used as ground truth must be exact to solver
        tolerance, since it forms the denominator of the efficiency metric.
\end{enumerate}

\section{Parameter identification}

\subsection{The five constraints}

There are five unknowns, $\{I_{ph}, I_0, R_s, R_{sh}, a\}$. Five independent
conditions are imposed, all evaluated at STC:

\begin{align}
\text{short circuit:} \quad & I(0) = \Isc \label{eq:c1}\\
\text{open circuit:} \quad & I(\Voc) = 0 \label{eq:c2}\\
\text{maximum power point:} \quad & I(\Vmp) = \Imp \label{eq:c3}\\
\text{stationarity at the MPP:} \quad & \Imp + \Vmp \left.\frac{\mathrm{d}I}{\mathrm{d}V}\right|_{\Vmp} = 0 \label{eq:c4}\\
\text{temperature coefficient:} \quad & \left.\frac{\mathrm{d}\Voc}{\mathrm{d}T}\right|_{T_{ref}} = \beta_{Voc} \label{eq:c5}
\end{align}

Conditions~\eqref{eq:c1} to~\eqref{eq:c3} place the curve through three known
points. Condition~\eqref{eq:c4} is the crucial one: it forces the model to have
the correct \emph{slope} at the MPP, which is what fixes the curvature and hence
$r_d$. Condition~\eqref{eq:c5} closes the system using the published temperature
coefficient, which is the De Soto closure and uses information that simpler
identification schemes discard.

\subsection{Why not fix the ideality factor}

A common alternative, due to Villalva et al.~\cite{villalva}, fixes $a$ at an
assumed value and iterates on $R_s$ until the modelled peak power matches the
datasheet. That method is robust and widely used, but it leaves one parameter to
the modeller's judgement and does not use $\beta_{Voc}$ at all. Solving all five
equations simultaneously removes that arbitrariness, and the resulting ideality
factor becomes an \emph{output} that can be sanity-checked against physics: a
value near unity is expected for a good crystalline silicon cell.

\subsection{Numerical formulation}

The system is solved with a trust-region method on the residual vector. Three
implementation choices make it reliable:

\begin{itemize}[leftmargin=1.4em]
  \item \textbf{Log parameterisation.} The solver works with
        $\left[I_{ph},\ \ln I_0,\ \ln R_s,\ \ln R_{sh},\ n\right]$, which keeps
        every physical quantity positive by construction and conditions the
        problem, since $I_0$ spans ten orders of magnitude.
  \item \textbf{Residual scaling.} The four current residuals are divided by
        $\Isc$ and the temperature residual by $|\beta_{Voc}|$, so all five
        entries are of comparable size and the trust region is meaningful.
  \item \textbf{Multi-start.} The solve is repeated from six ideality-factor
        seeds between \num{0.8} and \num{2.0}, and the best converged result is
        kept, so the answer does not depend on a lucky initial guess.
\end{itemize}

\subsection{Results}

Table~\ref{tab:pvfit} shows the identified parameters and
Table~\ref{tab:pvvalid} the resulting agreement with the datasheet. All five
constraints are satisfied to machine precision, and the extracted ideality
factor of \num{1.0036} is physically sensible for a multicrystalline silicon
cell.

\begin{table}[htbp]
\centering
\caption{Identified reference parameters for the Kyocera KC200GT, module level, at standard test conditions.}
\label{tab:pvfit}
\begin{tabular}{@{}llll@{}}
\toprule
\textbf{Parameter} & \textbf{Symbol} & \textbf{Value} & \textbf{Unit} \\
\midrule
Photocurrent            & $I_{ph,ref}$  & \num{8.22713}    & \si{\ampere} \\
Saturation current      & $I_{0,ref}$   & \num{4.3964e-10} & \si{\ampere} \\
Series resistance       & $R_{s,ref}$   & \num{0.33501}    & \si{\ohm} \\
Shunt resistance        & $R_{sh,ref}$  & \num{160.6015}   & \si{\ohm} \\
Modified ideality factor& $a_{ref}$     & \num{1.39246}    & \si{\volt} \\
Diode ideality factor   & $n$           & \num{1.0036}     & -- \\
\bottomrule
\end{tabular}
\end{table}

\begin{table}[htbp]
\centering
\caption{Model against datasheet at standard test conditions. The residual is at the level of double-precision round-off.}
\label{tab:pvvalid}
\begin{tabular}{@{}lrrr@{}}
\toprule
\textbf{Quantity} & \textbf{Datasheet} & \textbf{Model} & \textbf{Error (\si{\percent})} \\
\midrule
$\Voc$ (\si{\volt})            & 32.90    & 32.90    & \num{0} \\
$\Isc$ (\si{\ampere})          & 8.21     & 8.21     & \num{0} \\
$\Vmp$ (\si{\volt})            & 26.30    & 26.30    & \num{-1.4e-14} \\
$\Imp$ (\si{\ampere})          & 7.61     & 7.61     & \num{0} \\
$\Pmp$ (\si{\watt})            & 200.14   & 200.14   & \num{0} \\
$\beta_{Voc}$ (\si{\volt\per\kelvin}) & -0.12301 & -0.12301 & \num{1.1e-14} \\
\bottomrule
\end{tabular}
\end{table}

The same procedure applied to a SunPower SPR-305E-WHT-D, a 96-cell
monocrystalline module of quite different construction, converges equally well
with a worst error of \SI{7e-11}{\percent}, which indicates that the method is
not tuned to one datasheet.

\section{Solving the implicit equation}
\label{sec:solver}

Because~\eqref{eq:singlediode} cannot be inverted in closed form, terminal
current is obtained by Newton iteration on the residual

\begin{equation}
f(I) = I_{ph} - I_0\left[\exp\!\left(\tfrac{V + I R_s}{a}\right) - 1\right] - \frac{V + I R_s}{R_{sh}} - I .
\end{equation}

Three properties make this safe. First, $f$ is strictly decreasing in $I$, so
the root is unique. Second, starting from $I = 0$ the residual is positive
whenever $V < \Voc$, and for a concave decreasing function Newton's method
converges monotonically from that side. Third, the exponent is clamped and a
backtracking line search halves the step whenever the residual magnitude would
grow, which covers the transient excursions beyond $\Voc$ that a controller can
briefly demand.

In simulation the previous sample's current is supplied as the initial guess.
This warm start reduces the iteration count to two or three and is what makes
second-scale closed-loop simulation practical.

\section{Derived quantities}

The slope is evaluated from the analytic expression~\eqref{eq:didv} rather than
by differencing. Verification against a central difference over the range
$0.05\,\Voc$ to $0.97\,\Voc$ agrees to a relative error below \num{1e-5}, which
is the accuracy floor of the difference itself, not of the analytic formula.

The maximum power point is found by scanning \num{400} points to bracket the
peak and then root-solving the analytic $\mathrm{d}P/\mathrm{d}V$ inside that
bracket. The stationarity residual at the returned point is below
\SI{1e-6}{\watt\per\volt}.

\section{Behaviour of the model}

\begin{figure}[htbp]
\centering
\includegraphics[width=\textwidth]{s01_pv_characteristics.png}
\caption[Photovoltaic array characteristics]{Characteristics of the 6 series by 2 parallel KC200GT array. Top left, I--V curves at \SI{25}{\celsius}; top right, the corresponding P--V curves with a single interior maximum; bottom left, the effect of cell temperature at \SI{1000}{\watt\per\square\metre}; bottom right, dynamic resistance $r_d$ on a logarithmic axis, which is the quantity that sets the plant gain.}
\label{fig:pvchar}
\end{figure}

Figure~\ref{fig:pvchar} shows the identified model in use.
Table~\ref{tab:pvops} lists the array operating points.

\begin{table}[htbp]
\centering
\caption{Array operating points at \SI{25}{\celsius}, 6 series by 2 parallel.}
\label{tab:pvops}
\begin{tabular}{@{}rrrrrrrr@{}}
\toprule
$G$ & $\Vmp$ & $\Imp$ & $\Pmp$ & $\Voc$ & $\Isc$ & $FF$ & $r_d$ at MPP \\
(\si{\watt\per\square\metre}) & (\si{\volt}) & (\si{\ampere}) & (\si{\watt}) & (\si{\volt}) & (\si{\ampere}) & & (\si{\ohm}) \\
\midrule
200  & 156.02 & 3.061  & 477.6  & 183.97 & 3.290  & 0.789 & 50.97 \\
400  & 158.75 & 6.117  & 971.1  & 189.75 & 6.576  & 0.778 & 25.95 \\
600  & 159.21 & 9.163  & 1458.8 & 193.14 & 9.860  & 0.766 & 17.38 \\
800  & 158.75 & 12.198 & 1936.4 & 195.54 & 13.141 & 0.754 & 13.02 \\
1000 & 157.80 & 15.220 & 2401.7 & 197.40 & 16.420 & 0.741 & 10.37 \\
\bottomrule
\end{tabular}
\end{table}

Three observations from these results shape everything that follows.

\subsection{The MPP voltage is remarkably stable against irradiance}

Across a fivefold change in irradiance, $\Vmp$ moves only from
\SI{156.0}{\volt} to \SI{159.2}{\volt}, about \SI{2}{\percent}. This is why
regulating PV \emph{voltage} is a sensible inner objective: the target barely
moves, so the outer MPPT layer only has to make small corrections. It is also
the basis of the crude ``fractional open-circuit voltage'' MPPT methods, which
simply hold $\vpv \approx 0.78\,\Voc$.

\subsection{Temperature moves the MPP far more than irradiance does}

The bottom-left panel of Figure~\ref{fig:pvchar} shows $\Vmp$ falling
substantially as cell temperature rises, because $\Voc$ falls at
\SI{-0.123}{\volt\per\kelvin} per module, that is \SI{-0.738}{\volt\per\kelvin}
for the six-module string. Over a realistic \SI{40}{\kelvin} swing that is
nearly \SI{30}{\volt}. Temperature, not irradiance, is the dominant influence on
where the MPP sits, even though irradiance dominates how much power is there.

\subsection{The dynamic resistance varies enormously}

At the MPP, $r_d$ ranges from \SI{10.4}{\ohm} at full sun to \SI{51.0}{\ohm} at
\SI{200}{\watt\per\square\metre}. Along the curve at fixed irradiance the swing
is far larger still, from under \SI{2}{\ohm} near $\Voc$ to nearly
\SI{500}{\ohm} on the current-source branch. Chapter~\ref{ch:converter} shows
that this variation maps directly onto the damping of the converter's dominant
mode, which is the central difficulty of the control problem.

\chapter{The Converter and its Small-Signal Model}
\label{ch:converter}

\section{Topology and the choice of controlled variable}

The power stage is a boost converter with the array on its input and a stiff
\SI{400}{\volt} DC bus on its output, drawn in Figure~\ref{fig:boostckt}. The
bus is treated as stiff because in a real string inverter the downstream
grid-side stage regulates it; from the point of view of the MPPT converter it is
a voltage sink.

\begin{figure}[htbp]
\centering
\begin{circuitikz}[scale=1.0, transform shape]
\draw (0,0) to[isource, l_=$\ipv(\vpv)$, invert] (0,3.2);
\draw (0,3.2) -- (1.6,3.2);
\draw (1.6,3.2) to[C, l_=$C_{\mathrm{pv}}$] (1.6,0);
\draw (1.6,3.2) to[L, l=$L$, i>^=$i_L$] (4.2,3.2)
      to[R, l=$R_L$] (5.8,3.2);
\draw (5.8,3.2) to[Tnmos, n=mos] (5.8,0.6);
\draw (5.8,0.6) -- (5.8,0);
\draw (5.8,3.2) to[D, l=$D$] (8.2,3.2);
\draw (8.2,3.2) to[C, l_=$C_o$] (8.2,0);
\draw (8.2,3.2) -- (9.8,3.2) to[battery1, l=$V_{dc}$] (9.8,0);
\draw (0,0) -- (9.8,0);
\draw (0,0) node[ground]{};
\node at (0.55,1.6) {$\vpv$};
\draw[dashed] (-0.55,-0.5) rectangle (0.9,3.8);
\node[align=center, font=\small] at (0.2,4.2) {PV array};
\node[font=\small] at (mos.gate) [left=6pt] {$d$};
\end{circuitikz}
\caption[Boost converter power stage]{The boost converter power stage. The array is a nonlinear source characterised by $\ipv(\vpv)$; the duty ratio $d$ of the transistor is the single control input.}
\label{fig:boostckt}
\end{figure}

A design decision is embedded here that deserves emphasis. The converter has one
control input, the duty ratio, and several signals could be regulated with it:
the output current into the bus, the inductor current, or the PV terminal
voltage. This project regulates \textbf{PV terminal voltage} as the outer
objective. Two reasons:

\begin{enumerate}[leftmargin=1.6em]
  \item As Chapter~\ref{ch:pvmodel} showed, $\Vmp$ barely moves with irradiance,
        so a voltage reference is a naturally well-conditioned target.
  \item As shown in Section~\ref{sec:nozero}, the transfer function from duty
        ratio to PV voltage has no right half plane zero, whereas the
        transfer function from duty ratio to boost \emph{output} voltage
        famously does. Choosing the input-side variable avoids a fundamental
        bandwidth limitation.
\end{enumerate}

\section{Averaged model}

\subsection{Deriving the state equations}

Take as states the PV terminal voltage $\vpv$, which is the voltage on
$C_{\mathrm{pv}}$, and the inductor current $i_L$. Averaging the two switch
configurations over one switching period, weighted by $d$ and $1-d$, gives

\begin{align}
C_{\mathrm{pv}}\,\frac{\mathrm{d}\vpv}{\mathrm{d}t} &= \ipv(\vpv) - i_L \label{eq:av1}\\
L\,\frac{\mathrm{d}i_L}{\mathrm{d}t} &= \vpv - (1-d)\,V_{dc} - i_L R_L \label{eq:av2}
\end{align}

Equation~\eqref{eq:av1} is charge balance on the input capacitor: whatever the
array supplies and the inductor does not take, charges $C_{\mathrm{pv}}$.
Equation~\eqref{eq:av2} is the inductor voltage: the array voltage drives it
forward, the reflected bus voltage $(1-d)V_{dc}$ opposes it, and the winding
resistance dissipates.

The nonlinearity sits entirely in $\ipv(\vpv)$, the array characteristic.

\subsection{Steady state}

Setting the derivatives to zero gives $i_L = \ipv(\vpv)$, which simply says the
converter draws exactly what the array produces, and

\begin{equation}
D = 1 - \frac{\vpv - i_L R_L}{V_{dc}},
\label{eq:dutyexact}
\end{equation}

which is the ideal ratio~\eqref{eq:boostratio} corrected for the inductor
resistance drop.

\section{Linearisation}

\subsection{The state matrices}

Perturbing about an operating point $(V, I, D)$ and writing
$g = \mathrm{d}I/\mathrm{d}V$ for the array slope there,

\begin{equation}
A = \begin{bmatrix}
\dfrac{g}{C_{\mathrm{pv}}} & -\dfrac{1}{C_{\mathrm{pv}}}\\[2ex]
\dfrac{1}{L} & -\dfrac{R_L}{L}
\end{bmatrix},
\qquad
B = \begin{bmatrix} 0 \\[1ex] \dfrac{V_{dc}}{L}\end{bmatrix} .
\label{eq:AB}
\end{equation}

Three outputs are of interest, PV voltage, inductor current, and PV power:

\begin{equation}
C = \begin{bmatrix}
1 & 0\\
0 & 1\\
I + Vg & 0
\end{bmatrix} .
\label{eq:C}
\end{equation}

The third row deserves attention: the linearised gain from a state perturbation
to power is $I + Vg$, which by~\eqref{eq:dpdv} is exactly
$\mathrm{d}P/\mathrm{d}V$. Section~\ref{sec:vanishing} draws out the
consequence.

\subsection{Verification}

The matrices in~\eqref{eq:AB} were checked against a central-difference Jacobian
of the nonlinear right-hand side~\eqref{eq:av1}--\eqref{eq:av2}. At
\SI{700}{\watt\per\square\metre} and \SI{40}{\celsius} the relative discrepancy
is \num{7.4e-11} for $A$ and \num{3.9e-10} for $B$, and the residual at the
computed operating point is \SI{1.1e-11}{}. This check is part of the automated
test suite, so a future edit that breaks the analytic derivative fails a test
rather than silently corrupting every downstream result.

\section{The duty-to-voltage transfer function}
\label{sec:nozero}

Evaluating $G_{vd}(s) = C_1 (sI - A)^{-1} B$ gives

\begin{equation}
G_{vd}(s) = \frac{\vpv(s)}{d(s)}
= \frac{-\,\dfrac{V_{dc}}{L C_{\mathrm{pv}}}}
       {s^2 + \left(\dfrac{R_L}{L} - \dfrac{g}{C_{\mathrm{pv}}}\right)s
        + \dfrac{1 - g R_L}{L C_{\mathrm{pv}}}} .
\label{eq:Gvd}
\end{equation}

Several things can be read straight off this expression.

\subsection{No right half plane zero}

The numerator is a constant. There is no zero at all, let alone one in the right
half plane. This is the payoff of regulating the input-side voltage. The
non-minimum-phase behaviour of the boost converter arises in the duty-to-output
voltage path, where increasing duty momentarily \emph{reduces} output current
before the inductor current has time to build; that mechanism does not appear on
the input side.

\subsection{Sign of the gain}

The DC gain is $-V_{dc}/(1 - g R_L)$, which is negative because $g < 0$. Raising
the duty ratio lowers the PV voltage, as physical reasoning predicted. The
controller must therefore act on the error $\vpv - V^\ast$ rather than
$V^\ast - \vpv$, or equivalently carry a sign inversion.

\subsection{Natural frequency and damping}

Comparing~\eqref{eq:Gvd} with the standard second-order form,

\begin{align}
\omega_n &= \sqrt{\frac{1 - g R_L}{L C_{\mathrm{pv}}}} \;\approx\; \frac{1}{\sqrt{L C_{\mathrm{pv}}}}, \label{eq:wn}\\
2\zeta\omega_n &= \frac{R_L}{L} - \frac{g}{C_{\mathrm{pv}}} = \frac{R_L}{L} + \frac{1}{r_d C_{\mathrm{pv}}}. \label{eq:zeta}
\end{align}

This is the key structural result of the chapter. The natural frequency is set
by the passive components and barely depends on operating point. The damping,
however, contains the term $1/(r_d C_{\mathrm{pv}})$, so it is inversely
proportional to the array's dynamic resistance. Since $r_d$ varies by more than
two orders of magnitude along the I--V curve, so does the damping.

Physically, the array is acting as the damping element of the LC circuit. On the
voltage-source branch it is a low resistance and damps the resonance heavily; on
the current-source branch it is nearly an open circuit and leaves the LC
essentially undamped, with only the small winding resistance $R_L$ to dissipate.

\section{The operating envelope}
\label{sec:sizing}

\subsection{Component sizing}

Sizing followed from a \SI{20}{\percent} inductor current ripple target at the
STC maximum power point. From~\eqref{eq:iripple} with $\Vmp = \SI{157.8}{\volt}$,
$D = \num{0.6055}$ and $\Delta i_L = \SI{3.04}{\ampere}$, the required inductance
is \SI{1.57}{\milli\henry}, rounded to \SI{1.5}{\milli\henry}. The input
capacitor was then chosen for a PV voltage ripple below \SI{0.25}{\volt} using
$\Delta \vpv = \Delta i_L / (8 f_{sw} C_{\mathrm{pv}})$, giving
\SI{76}{\micro\farad}, rounded up to \SI{100}{\micro\farad}. The resulting LC
corner is \SI{411}{\hertz}, comfortably below $f_{sw}/10 = \SI{2}{\kilo\hertz}$
as averaging requires.

\subsection{Envelope in irradiance}

Table~\ref{tab:envG} sweeps irradiance with the array held at its maximum power
point. The last column confirms that the converter stays in continuous
conduction mode down to \SI{200}{\watt\per\square\metre}, where the mean current
is still \num{1.9} times half the ripple.

\begin{table}[htbp]
\centering
\caption{Small-signal properties at the maximum power point across irradiance, \SI{25}{\celsius}.}
\label{tab:envG}
\begin{tabular}{@{}rrrrrrrr@{}}
\toprule
$G$ & $\Vmp$ & $\Imp$ & $D$ & $r_d$ & $f_n$ & $\zeta$ & CCM ratio \\
(\si{\watt\per\square\metre}) & (\si{\volt}) & (\si{\ampere}) & & (\si{\ohm}) & (\si{\hertz}) & & \\
\midrule
200  & 156.02 & 3.06  & 0.611 & 50.97 & 411.5 & 0.0573 & 1.93 \\
400  & 158.75 & 6.12  & 0.605 & 25.95 & 412.1 & 0.0937 & 3.82 \\
600  & 159.21 & 9.16  & 0.605 & 17.38 & 412.7 & 0.1303 & 5.70 \\
800  & 158.75 & 12.20 & 0.608 & 13.02 & 413.3 & 0.1672 & 7.59 \\
1000 & 157.80 & 15.22 & 0.611 & 10.37 & 413.9 & 0.2047 & 9.47 \\
\bottomrule
\end{tabular}
\end{table}

\subsection{Envelope along the I--V curve}

Table~\ref{tab:envV} is the more revealing sweep. Holding irradiance at
\SI{1000}{\watt\per\square\metre} and moving the operating point along the
curve, the resonant frequency changes by under \SI{4}{\percent} while the
damping ratio changes by a factor of \num{44}.

\begin{table}[htbp]
\centering
\caption{Small-signal properties along the I--V curve at \SI{1000}{\watt\per\square\metre}. The resonance is essentially fixed; the damping is not.}
\label{tab:envV}
\begin{tabular}{@{}rrrrrr@{}}
\toprule
$V$ (\si{\volt}) & $V/\Voc$ & $r_d$ (\si{\ohm}) & $f_n$ (\si{\hertz}) & $\zeta$ & $\mathrm{d}P/\mathrm{d}V$ (\si{\watt\per\volt}) \\
\midrule
69.09  & 0.35 & 482.1 & 411.0 & 0.0234 & $+16.13$ \\
108.57 & 0.55 & 417.6 & 411.0 & 0.0240 & $+15.93$ \\
138.18 & 0.70 & 77.23 & 411.3 & 0.0444 & $+14.25$ \\
157.80 & 0.799 & 10.37 & 413.9 & 0.2047 & $\approx 0$ \\
173.71 & 0.88 & 3.077 & 420.8 & 0.6335 & $-44.40$ \\
187.53 & 0.95 & 1.832 & 427.4 & 1.0350 & $-96.41$ \\
\bottomrule
\end{tabular}
\end{table}

\begin{figure}[htbp]
\centering
\includegraphics[width=\textwidth]{s02_plant_frequency_response.png}
\caption[Small-signal plant frequency response]{Small-signal behaviour of the plant. Top left and top right, magnitude and phase of duty to PV voltage at six points along the I--V curve: the resonance peak grows dramatically as the operating point moves onto the current-source branch. Bottom left, duty to inductor current is almost invariant, which is what makes the inner loop easy. Bottom right, the DC gain from duty to PV power passes through zero exactly at the maximum power point and changes sign.}
\label{fig:plantbode}
\end{figure}

Figure~\ref{fig:plantbode} shows the same information as frequency responses.
The resonance peak of $G_{vd}$ reaches roughly \SI{27}{\decibel} above its low
frequency value on the current-source branch, and essentially disappears near
$\Voc$.

\section{Why power cannot be the feedback variable}
\label{sec:vanishing}

The bottom-right panel of Figure~\ref{fig:plantbode} makes concrete a point that
is easy to state and easy to overlook. Combining the third row of~\eqref{eq:C}
with~\eqref{eq:Gvd}, the DC gain from duty ratio to PV power is

\begin{equation}
\left.\frac{\partial P}{\partial d}\right|_{DC}
= \frac{\mathrm{d}P}{\mathrm{d}V}\cdot \frac{-V_{dc}}{1 - g R_L} .
\label{eq:powergain}
\end{equation}

At the maximum power point $\mathrm{d}P/\mathrm{d}V = 0$ by definition, so
\emph{the plant gain from the control input to the power output is exactly
zero there}, and it changes sign as the operating point crosses the peak. The
numerically computed value at the MPP is \SI{8.9e-15}{\watt}, that is, zero to
round-off.

Two consequences follow directly:

\begin{enumerate}[leftmargin=1.6em]
  \item \textbf{Power cannot be used as a conventional feedback error.} A linear
        controller closing a loop on power would have infinite required gain
        exactly where it needs to work, and would have the wrong sign on one
        side of the peak. No fixed linear controller can regulate to an
        extremum.
  \item \textbf{MPPT must be a fundamentally different kind of layer.} It must
        estimate the \emph{gradient} rather than regulate the value. Perturb and
        observe does this by finite differencing; incremental conductance by
        evaluating the analytic condition; extremum seeking by dither
        demodulation. Chapter~\ref{ch:mppt} treats all three as gradient
        estimators, which is what they are.
\end{enumerate}

This is the sense in which the project title is more than a framing device. The
control-theoretic view explains \emph{why} MPPT algorithms have the form they
do, rather than presenting them as a catalogue of heuristics.

\chapter{Control Design and Analysis}
\label{ch:control}

\section{The design problem}

Chapter~\ref{ch:converter} established the difficulty precisely. The plant
$G_{vd}$ is a second-order resonance whose frequency is essentially fixed at
about \SI{412}{\hertz} but whose damping ratio varies from \num{0.023} to
\num{1.035} depending only on where the array sits on its I--V curve. A
controller must hold adequate stability margins across that entire range, using
one fixed set of gains, because the operating point is dictated by the MPPT
layer and by the weather rather than by the designer.

Two candidate architectures are compared.

\begin{figure}[htbp]
\centering
\resizebox{\textwidth}{!}{%
\begin{tikzpicture}[
  scale=0.88, transform shape,
  font=\small,
  box/.style={draw, rounded corners=2pt, minimum height=9mm, minimum width=24mm, align=center},
  >={Stealth[length=2.2mm]}
]
\node[box] (mppt) {MPPT\\ \SI{50}{\hertz}};
\node[box, right=10mm of mppt] (vloop) {Voltage PI\\ \SI{250}{\hertz}};
\node[box, right=10mm of vloop] (iloop) {Current PI\\ \SI{2}{\kilo\hertz}};
\node[box, right=10mm of iloop] (pwm) {PWM\\ \SI{20}{\kilo\hertz}};
\node[box, right=10mm of pwm] (plant) {Converter\\ and array};

\draw[->] (mppt) -- node[above]{$V^\ast$} (vloop);
\draw[->] (vloop) -- node[above]{$i_L^\ast$} (iloop);
\draw[->] (iloop) -- node[above]{$d$} (pwm);
\draw[->] (pwm) -- (plant);

\coordinate (out) at ($(plant.east)+(9mm,0)$);
\draw[->] (plant.east) -- (out) node[right]{$\vpv,\ \ipv,\ i_L$};

\draw[->] ($(plant.east)+(5mm,0)$) |- ($(iloop.south)+(0,-7mm)$) -| (iloop.south)
  node[pos=0.25, below]{$i_L$};
\draw[->] ($(plant.east)+(7mm,0)$) |- ($(vloop.south)+(0,-14mm)$) -| (vloop.south)
  node[pos=0.25, below]{$\vpv$};
\draw[->] ($(plant.east)+(9mm,0)$) |- ($(mppt.south)+(0,-21mm)$) -| (mppt.south)
  node[pos=0.25, below]{$\vpv,\ \ipv$};
\end{tikzpicture}}
\caption[Cascaded control architecture]{The cascaded architecture. Each layer commands the reference of the layer to its left and runs roughly eight times slower than the layer to its right.}
\label{fig:cascade}
\end{figure}

\begin{description}[leftmargin=0pt, style=unboxed]
\item[Single-loop.] One proportional--integral controller acting on the PV
  voltage error and producing the duty ratio directly. This is the obvious first
  attempt and appears widely in the MPPT literature.
\item[Cascade.] An inner PI regulating inductor current, with an outer PI
  regulating PV voltage by commanding the inner current reference, as in
  Figure~\ref{fig:cascade}.
\end{description}

\section{The single-loop design and why it fails}

\subsection{Design}

The single-loop controller was designed conventionally: crossover placed at
\SI{130}{\hertz}, roughly a third of the LC resonance, with the PI zero a factor
of five below crossover at \SI{26}{\hertz}. Evaluating the plant magnitude at
crossover and solving $|L(j\omega_c)| = 1$ gives $K_p = \num{2.264e-3}$ and
$K_i = \num{0.370}$. Designed at the STC maximum power point, this is a
perfectly reasonable-looking design.

\subsection{The multiple-crossover problem}

The trouble appears immediately when the loop gain is examined over frequency
rather than at a single point. Because a PI contributes an approximately flat
gain above its zero, and because the plant magnitude is flat below its
resonance, the loop gain is nearly flat between the intended crossover and the
resonance. The resonance then lifts the magnitude by up to \SI{27}{\decibel},
pushing $|L|$ back above unity, and it must come down through unity a third
time on the far side, by which point the plant has contributed its full
\SI{-180}{\degree}.

The measured result is three gain crossovers, at \SI{56}{\hertz},
\SI{130}{\hertz} and \SI{542}{\hertz}, with a peak loop gain of
\SI{43.4}{\decibel}. When a loop crosses unity more than once, quoting the
margin at the intended crossover is meaningless; the relevant number is the
worst margin over all crossings.

\subsection{Collapse of the phase margin}

Table~\ref{tab:singleloop} gives the worst-case phase margin along the I--V
curve. On the current-source branch, where the damping ratio is \num{0.023}, the
phase margin falls to \SI{1.5}{\degree}. The loop is nominally stable, but a
\SI{1.5}{\degree} margin is no margin at all: any unmodelled delay, any sensor
filter, any component tolerance destroys it, and even nominally it would ring
almost undamped for a very long time.

\begin{table}[htbp]
\centering
\caption{Worst-case phase margin along the I--V curve at \SI{1000}{\watt\per\square\metre}. The single-loop controller degrades exactly where the array behaves as a current source.}
\label{tab:singleloop}
\begin{tabular}{@{}rrrrrr@{}}
\toprule
$V$ (\si{\volt}) & $r_d$ (\si{\ohm}) & $\zeta$ & \multicolumn{2}{c}{Single-loop PI} & Cascade outer \\
\cmidrule(lr){4-5}\cmidrule(lr){6-6}
 & & & PM (\si{\degree}) & crossings & PM (\si{\degree}) \\
\midrule
69.09  & 482.1 & 0.0234 & 1.46  & 3 & 85.3 \\
88.83  & 475.8 & 0.0234 & 1.47  & 3 & 85.3 \\
108.57 & 417.6 & 0.0240 & 1.57  & 3 & 85.4 \\
128.31 & 182.9 & 0.0299 & 2.61  & 3 & 86.4 \\
138.18 & 77.23 & 0.0444 & 5.14  & 3 & 88.8 \\
148.05 & 27.94 & 0.0884 & 12.91 & 3 & 95.9 \\
157.92 & 10.25 & 0.2068 & 34.50 & 3 & 114.2 \\
167.79 & 4.458 & 0.4463 & 88.62 & 3 & 122.7 \\
177.66 & 2.539 & 0.7600 & 139.8 & 1 & 110.6 \\
187.53 & 1.832 & 1.0350 & 135.4 & 1 & 105.3 \\
\bottomrule
\end{tabular}
\end{table}

This is not a matter of retuning. The problem is structural: a PI has no
mechanism for damping a resonance it cannot see, and the required gain reduction
to keep $|L|$ below unity through the resonance would push crossover down to a
few hertz, making the loop uselessly slow.

\section{The cascaded design}

\subsection{Why an inner current loop fixes the problem}

Consider the transfer function from duty ratio to inductor current,

\begin{equation}
G_{id}(s) = \frac{i_L(s)}{d(s)}
= \frac{V_{dc}}{L}\cdot
  \frac{s - g/C_{\mathrm{pv}}}
       {s^2 + \left(\tfrac{R_L}{L} - \tfrac{g}{C_{\mathrm{pv}}}\right)s + \tfrac{1-gR_L}{L C_{\mathrm{pv}}}} .
\label{eq:Gid}
\end{equation}

The zero sits at $s = g/C_{\mathrm{pv}}$, which is in the left half plane
because $g < 0$, so this path is minimum phase. More importantly, at frequencies
well above the resonance the expression collapses to

\begin{equation}
G_{id}(s) \;\approx\; \frac{V_{dc}}{L s},
\label{eq:Gidhf}
\end{equation}

a pure integrator whose gain depends only on $V_{dc}$ and $L$ and \emph{not at
all} on the operating point. This is visible directly in the bottom-left panel
of Figure~\ref{fig:plantbode}, where the six curves are indistinguishable above
a few hundred hertz.

So an inner loop closed at \SI{2}{\kilo\hertz}, comfortably above the
\SI{412}{\hertz} resonance, sees an operating-point-independent plant. Its gains
follow immediately from~\eqref{eq:Gidhf}: for unity loop gain at
$\omega_{ci} = 2\pi \times \SI{2000}{\hertz}$,

\begin{equation}
K_{p,i} = \frac{L\,\omega_{ci}}{V_{dc}} = \num{0.04712},
\qquad
K_{i,i} = K_{p,i}\frac{\omega_{ci}}{10} = \num{59.2},
\end{equation}

placing the PI zero a decade below crossover. Because the plant contributes
\SI{-90}{\degree} there and the PI adds only a few degrees more, the phase
margin is close to \SI{85}{\degree} by construction.

\subsection{The outer loop sees a first-order plant}

With the inner loop fast and accurate, $i_L \approx i_L^\ast$, and the input
capacitor equation~\eqref{eq:av1} linearises to

\begin{equation}
\frac{\tilde v_{\mathrm{pv}}(s)}{\tilde i_L^\ast(s)}
= \frac{-1}{C_{\mathrm{pv}} s - g}
= \frac{-1}{C_{\mathrm{pv}} s + 1/r_d} .
\label{eq:outerplant}
\end{equation}

The resonance is gone. What remains is a single real pole at
$1/(r_d C_{\mathrm{pv}})$. Its DC gain is $r_d$, which still varies widely, but
at the outer crossover of \SI{250}{\hertz} the capacitive term dominates
whenever $\omega \gg 1/(r_d C_{\mathrm{pv}})$, so the plant looks like
$1/(C_{\mathrm{pv}} s)$ and the required proportional gain is again
operating-point independent:

\begin{equation}
K_{p,v} = C_{\mathrm{pv}}\,\omega_{cv} = \num{0.15708},
\qquad
K_{i,v} = K_{p,v}\frac{\omega_{cv}}{5} = \num{49.35}.
\end{equation}

Physically, the inner current loop turns the inductor into a controlled current
source, and a controlled current source charging a capacitor is about the
easiest plant in control engineering. That is the whole reason the cascade
works.

\subsection{Verification on the full-order model}

The gains above were derived from simplified models: the high-frequency
approximation~\eqref{eq:Gidhf} for the inner loop and the ideal-inner-loop
approximation~\eqref{eq:outerplant} for the outer. Design on a simplified model
is only legitimate if the result is then checked on the full one.

Both loops were therefore re-analysed on the exact second-order plant, with the
inner loop closed exactly rather than assumed ideal:

\begin{align}
L_i(s) &= C_i(s)\,G_{id}(s), \\
L_v(s) &= -C_v(s)\cdot G_{vd}(s)\,\frac{C_i(s)}{1 + C_i(s) G_{id}(s)} .
\end{align}

Results are in the last column of Table~\ref{tab:singleloop} and in
Table~\ref{tab:cascadeG}. The inner loop margin is between \SI{85.0}{\degree}
and \SI{85.8}{\degree} everywhere, confirming operating-point independence. The
outer loop margin is between \SI{85.3}{\degree} and \SI{122.7}{\degree}, always
far above the \SI{45}{\degree} target.

\begin{table}[htbp]
\centering
\caption{Cascade margins at the maximum power point across irradiance, verified on the full-order plant.}
\label{tab:cascadeG}
\begin{tabular}{@{}rrrrrr@{}}
\toprule
$G$ (\si{\watt\per\square\metre}) & $r_d$ (\si{\ohm}) & Inner PM (\si{\degree}) & Inner $f_c$ (\si{\hertz}) & Outer PM (\si{\degree}) & Outer $f_c$ (\si{\hertz}) \\
\midrule
200  & 50.97 & 85.02 & 2090 & 90.9  & 215.7 \\
400  & 25.95 & 85.06 & 2090 & 96.8  & 205.0 \\
600  & 17.38 & 85.09 & 2090 & 102.6 & 191.4 \\
800  & 13.02 & 85.12 & 2090 & 108.4 & 174.7 \\
1000 & 10.37 & 85.16 & 2089 & 113.9 & 155.4 \\
\bottomrule
\end{tabular}
\end{table}

\begin{figure}[htbp]
\centering
\includegraphics[width=\textwidth]{s03_control_design.png}
\caption[Control design comparison]{Comparison of the two architectures. Top left, single-loop gain magnitude at five operating points: the resonance forces the curve back above \SI{0}{\decibel}. Top right, the cascade outer loop crosses unity once and cleanly. Bottom left, worst-case phase margin against operating voltage, the central result of the chapter. Bottom right, closed-loop reference step responses of the cascade at five operating points.}
\label{fig:controldesign}
\end{figure}

\subsection{Bandwidth variation and what it costs}

One genuine operating-point dependence remains in the cascade: the outer
crossover falls from about \SI{223}{\hertz} on the current-source branch to
\SI{15}{\hertz} near $\Voc$. This happens because the plant DC gain $r_d$ has
become so small there that the loop gain drops below unity at a lower frequency.
The bottom-right panel of Figure~\ref{fig:controldesign} shows the
consequence directly: the settling time varies from a few milliseconds to about
\SI{50}{\milli\second} across the range.

This variation is benign, because the slow region is far from the maximum power
point where the tracker actually operates, and because the phase margin remains
above \SI{100}{\degree} throughout. If uniform bandwidth were required, the
outer proportional gain could be scheduled on the measured $r_d$, which is
directly available as $-\Delta \vpv / \Delta \ipv$ from the MPPT perturbations.
That refinement was not needed here and is noted as future work.

\subsection{Measured step response}

The design predicts a \SI{250}{\hertz} outer loop with a large phase margin,
therefore a well-damped response settling in roughly \SI{10}{\milli\second}.
Simulating the full nonlinear averaged plant with the discrete controllers, a
\SI{140}{\volt} to \SI{165}{\volt} reference step settles to within
\SI{2}{\percent} in \SI{11.9}{\milli\second} with no overshoot, and the
steady-state error is zero to solver precision. Measured against the analytic
I--V curve, the settled operating point matches to better than
\SI{0.05}{\percent} in both voltage and current.

\section{Consequences for the MPPT layer}

The voltage loop settling time of roughly \SI{12}{\milli\second} is not merely a
performance figure; it is a hard constraint on the layer above. A perturb and
observe tracker must wait for the voltage loop to settle before it measures the
power resulting from its perturbation, otherwise it measures a transient and
draws the wrong conclusion. That is why the MPPT sample period was set to
\SI{20}{\milli\second}, comfortably longer than the settling time.

This is timescale separation appearing as a design rule rather than as an
abstraction, and it is the reason a control-theoretic treatment of MPPT is
worth the effort: the number \SI{20}{\milli\second} is derived, not guessed.

\chapter{Maximum Power Point Tracking Algorithms}
\label{ch:mppt}

\section{All three algorithms are gradient estimators}

Section~\ref{sec:vanishing} established that the plant gain from control input
to power vanishes at the optimum, so no fixed linear controller can regulate to
it. What is needed instead is an estimate of the \emph{sign}, and ideally the
magnitude, of $\mathrm{d}P/\mathrm{d}V$. The three algorithms implemented here
differ only in how they obtain that estimate.

\begin{table}[htbp]
\centering
\caption{The three trackers viewed as gradient estimators.}
\label{tab:estimators}
\begin{tabular}{@{}p{0.22\textwidth}p{0.40\textwidth}p{0.28\textwidth}@{}}
\toprule
\textbf{Algorithm} & \textbf{How it estimates the gradient} & \textbf{Cost} \\
\midrule
Perturb and observe & Sign of $\Delta P$ between successive deliberate perturbations & Permanent limit cycle of $2\times$ the step \\
Incremental conductance & Direct evaluation of $\mathrm{d}I/\mathrm{d}V + I/V$ from successive samples & Same limit cycle, plus sensitivity to noise in a ratio \\
Extremum seeking & Demodulation of the response to a sinusoidal dither & Permanent dither of amplitude $A$ \\
\bottomrule
\end{tabular}
\end{table}

Every one of them pays a permanent steady-state cost, because every one of them
must keep exciting the plant to keep learning the gradient. That trade is
unavoidable and is quantified in Section~\ref{sec:tradeoff}.

\section{Perturb and observe}

\subsection{Principle}

The \term{perturb and observe} (P\&O) algorithm holds a direction $s \in \{+1,
-1\}$ and a fixed step $\Delta V$. Every sample period it moves the voltage
reference by $s\,\Delta V$, measures the resulting power, and reverses direction
if the power fell:

\begin{equation}
s_{k} = \begin{cases} -s_{k-1}, & P_k < P_{k-1} \\ \phantom{-}s_{k-1}, & \text{otherwise}\end{cases}
\qquad
V^\ast_{k+1} = V^\ast_k + s_k \Delta V .
\end{equation}

\subsection{The limit cycle}

At convergence the algorithm cannot stop; it has no way to know it has arrived.
It oscillates between three reference values, giving a steady-state limit cycle
of peak-to-peak amplitude exactly $2\Delta V$. Simulation confirms this
precisely: measured peak-to-peak values were \SI{1.00}{\volt},
\SI{1.99}{\volt}, \SI{2.99}{\volt}, \SI{3.98}{\volt}, \SI{5.96}{\volt},
\SI{7.98}{\volt} and \SI{11.95}{\volt} for steps of \SI{0.5}{\volt} to
\SI{6}{\volt}. This exactness is a useful validation that the voltage loop is
tracking the reference faithfully.

\subsection{The sampling constraint}

P\&O attributes the observed power change entirely to its own perturbation. Two
conditions must hold for that attribution to be valid:

\begin{enumerate}[leftmargin=1.6em]
  \item The voltage loop must have settled, or the measurement captures a
        transient. Chapter~\ref{ch:control} measured \SI{11.9}{\milli\second}
        settling, hence the \SI{20}{\milli\second} sample period.
  \item The irradiance must not have changed appreciably during the sample
        period. When it does, the attribution fails, which is the drift problem
        of Section~\ref{sec:drift}.
\end{enumerate}

\section{Incremental conductance}

\subsection{Principle}

\term{Incremental conductance} (INC)~\cite{hussein} rewrites the optimality
condition. Since $\mathrm{d}P/\mathrm{d}V = I + V\,\mathrm{d}I/\mathrm{d}V$,
setting it to zero gives

\begin{equation}
\frac{\mathrm{d}I}{\mathrm{d}V} = -\frac{I}{V},
\end{equation}

that is, the \emph{incremental} conductance equals the negative of the
\emph{instantaneous} conductance. The algorithm approximates the derivative by
$\Delta I / \Delta V$ from consecutive samples and moves the reference according
to the sign of the residual

\begin{equation}
r = \frac{\Delta I}{\Delta V} + \frac{I}{V} = \frac{1}{V}\frac{\mathrm{d}P}{\mathrm{d}V} .
\end{equation}

Because $V > 0$ always, $\mathrm{sign}(r) = \mathrm{sign}(\mathrm{d}P/\mathrm{d}V)$,
so the update rule is to increase the reference when $r > 0$ and decrease it
when $r < 0$. A tolerance band around zero is used to avoid chattering on noise.

\subsection{The case $\Delta V = 0$}

If the voltage has not changed between samples, the ratio is undefined. The
standard treatment, implemented here, examines $\Delta I$ instead: a positive
$\Delta I$ at constant voltage means the irradiance has risen, so the maximum
power point has moved to slightly higher voltage and the reference should
increase.

\subsection{Relationship to P\&O}

INC is often presented as fundamentally superior to P\&O. In the results of this
project the two are essentially identical under constant and slowly varying
irradiance, both reaching \SI{99.95}{\percent} at \SI{25}{\watt\per\square\metre\per\second},
and both showing the same $2\Delta V$ limit cycle. This should not be
surprising: with the same fixed step and the same sample period they are two
ways of computing the sign of the same quantity. INC's theoretical advantage,
that it evaluates a condition rather than differencing power, only materialises
when the two differ, which is under fast transients and in the presence of
measurement noise.

\section{Extremum seeking control}

\subsection{Structure}

The \term{extremum seeking} (ESC) tracker adds a sinusoidal dither to the
reference and extracts the gradient by correlation:

\begin{align}
V^\ast(t) &= \hat V(t) + A\sin(\omega t), \\
\xi(t) &= P(t) - \overline{P}(t) \quad\text{(washout high-pass)}, \\
\hat\gamma(t) &= \mathrm{LPF}\big[\xi(t)\sin(\omega t)\big] \;\approx\; \tfrac{A}{2}\frac{\mathrm{d}P}{\mathrm{d}V}, \\
\dot{\hat V}(t) &= K\,\hat\gamma(t) .
\end{align}

The washout filter removes the large DC component of power so that the
demodulator sees only the dither-induced variation. The low-pass filter after
demodulation removes the second-harmonic term. The integrator then performs
gradient ascent on the estimate.

\subsection{Parameter selection}

The dither frequency must satisfy competing constraints. It must be low enough
that the voltage control loop can actually follow the dither, so well below the
outer loop bandwidth of \SI{155}{\hertz} to \SI{223}{\hertz}. It must be high
enough that the dither is separable from genuine irradiance variation.
\SI{25}{\hertz} satisfies both. The washout and low-pass corners were set at
\SI{5}{\hertz}, a factor of five below the dither, and the amplitude at
\SI{3}{\volt}.

The adaptation gain $K$ was chosen by sweep, shown in Table~\ref{tab:escgain}.
The behaviour is exactly what averaging theory predicts: performance improves
with gain up to a point, then the estimator loop itself becomes unstable and the
steady-state excursion grows abruptly.

\begin{table}[htbp]
\centering
\caption{Extremum seeking gain sweep at \SI{1000}{\watt\per\square\metre}, starting from \SI{120}{\volt}. The true $\Vmp$ is \SI{157.80}{\volt}.}
\label{tab:escgain}
\begin{tabular}{@{}rrrr@{}}
\toprule
Gain $K$ & Static $\eta$ (\si{\percent}) & Estimate $\hat V$ (\si{\volt}) & Excursion (\si{\volt}) \\
\midrule
5   & 98.51 & 157.75 & 5.46 \\
20  & \textbf{99.21} & 157.75 & 5.51 \\
50  & 99.07 & 157.76 & 5.64 \\
100 & 98.62 & 157.79 & 5.89 \\
200 & 97.34 & 156.55 & 27.02 \\
400 & 97.67 & 161.31 & 27.06 \\
\bottomrule
\end{tabular}
\end{table}

$K = 20$ was adopted. Note that the internal estimate $\hat V$ converges to
\SI{157.75}{\volt} against a true value of \SI{157.80}{\volt}, an error of
\SI{0.03}{\percent}. ESC is the most \emph{accurate} of the three trackers; its
efficiency penalty comes entirely from the deliberate dither, not from bias.

\section{The fundamental trade-off}
\label{sec:tradeoff}

Every hill-climbing tracker faces the same trade: a larger perturbation acquires
the peak faster but oscillates further from it in steady state. This was
quantified by sweeping the P\&O step from \SI{0.5}{\volt} to \SI{6}{\volt},
measuring acquisition time from a \SI{120}{\volt} start and steady-state
efficiency separately, so that the two effects do not contaminate each other.

\begin{table}[htbp]
\centering
\caption{The P\&O speed against accuracy frontier at \SI{1000}{\watt\per\square\metre}.}
\label{tab:tradeoff}
\begin{tabular}{@{}rrrr@{}}
\toprule
Step (\si{\volt}) & Steady $\eta$ (\si{\percent}) & Limit cycle (\si{\volt}) & Acquisition (\si{\second}) \\
\midrule
0.5 & 99.9947 & 1.00  & 1.481 \\
1.0 & 99.9838 & 1.99  & 0.721 \\
1.5 & 99.9645 & 2.99  & 0.462 \\
2.0 & 99.9386 & 3.98  & 0.341 \\
3.0 & 99.7995 & 5.96  & 0.221 \\
4.0 & 99.6920 & 7.98  & 0.144 \\
6.0 & 99.4125 & 11.95 & 0.082 \\
\bottomrule
\end{tabular}
\end{table}

The relationship is cleanly monotone in both directions: an eighteen-fold
reduction in acquisition time costs \SI{0.58}{\percent} of steady-state
efficiency. A step of \SI{1.5}{\volt}, roughly \SI{1}{\percent} of $\Vmp$, was
adopted as a reasonable compromise and used for all comparisons.

\begin{figure}[htbp]
\centering
\includegraphics[width=\textwidth]{s04_mppt_comparison.png}
\caption[MPPT algorithm comparison]{Algorithm comparison through a \SI{1000}{} to \SI{500}{\watt\per\square\metre} irradiance step. Top left, voltage tracking; top right, power recovery; bottom left, steady-state behaviour showing the limit cycles and the extremum seeking dither; bottom right, the P\&O speed against accuracy frontier with each point labelled by its step size.}
\label{fig:mpptcompare}
\end{figure}

\section{Comparison under a step change}

Table~\ref{tab:stepcompare} compares all four tracker configurations through an
instantaneous irradiance step, the classic test.

\begin{table}[htbp]
\centering
\caption{Performance through a \SI{1000}{} to \SI{500}{\watt\per\square\metre} step.}
\label{tab:stepcompare}
\begin{tabular}{@{}lrrrr@{}}
\toprule
Tracker & $\eta$ (\si{\percent}) & RMS error (\si{\volt}) & Ripple (\si{\volt}) & Worst deficit (\si{\watt}) \\
\midrule
P\&O, fixed step        & 99.944 & 1.56 & 3.00 & 182.7 \\
P\&O, variable step     & \textbf{99.956} & 1.39 & \textbf{0.50} & \textbf{154.6} \\
Incremental conductance & 99.944 & 1.56 & 3.00 & 182.7 \\
Extremum seeking        & 99.551 & 4.88 & 5.95 & 194.6 \\
\bottomrule
\end{tabular}
\end{table}

Under this test the variable-step P\&O looks best: it shrinks its step near the
peak, achieving a \SI{0.5}{\volt} limit cycle, and enlarges it during the
transient. Extremum seeking looks worst, penalised by its \SI{6}{\volt}
peak-to-peak dither.

Chapter~\ref{ch:benchmark} shows that this ranking inverts completely under
irradiance \emph{ramps}, which is why a step test alone is an inadequate
benchmark and why the EN 50530 ramp profiles exist.

\section{The drift problem}
\label{sec:drift}

An instantaneous step is, paradoxically, an easy test for P\&O. The irradiance
is constant both before and after, so the power change the algorithm observes is
genuinely caused by its own perturbation except in the single sample that
straddles the step.

A \emph{ramp} is much harder. Suppose irradiance is rising at
\SI{100}{\watt\per\square\metre\per\second} and the sample period is
\SI{20}{\milli\second}. Between samples the irradiance rises by
\SI{2}{\watt\per\square\metre}, which at full sun adds roughly \SI{5}{\watt} to
the available power. A \SI{1.5}{\volt} perturbation near the flat peak changes
power by far less than that. So $\Delta P > 0$ regardless of which way the
algorithm perturbed, it concludes its last move was correct, and it keeps going
in that direction. If that direction happens to be wrong, it marches steadily
away from the peak. This is the well-documented \term{drift} or
\term{confusion} phenomenon~\cite{femia}.

Whether it occurs depends on the phase of the perturbation when the ramp starts,
which is why it is intermittent and easy to miss in a single simulation. The
next chapter observes it directly, and shows that a naive variable-step rule,
which was the best performer under the step test, makes it dramatically worse.

\chapter{Dynamic Benchmarking}
\label{ch:benchmark}

\section{Why a standard test is needed}

Chapter~\ref{ch:mppt} ended with four trackers separated by less than
\SI{0.5}{\percent} on a step test, in an order that turns out not to survive a
more demanding test. This is the normal situation in MPPT work: results depend
heavily on the test signal, so comparisons are only meaningful against an agreed
one.

The European standard EN 50530~\cite{en50530} defines such a test. It specifies
trapezoidal irradiance profiles, ramping linearly between two levels, dwelling,
ramping back, and repeating, with prescribed ramp rates. Two classes are
defined, a low-to-medium class between \SI{10}{\percent} and \SI{50}{\percent}
of nominal irradiance with gentle slopes, and a medium-to-high class between
\SI{30}{\percent} and \SI{100}{\percent} with steeper slopes. The figure of
merit is the dynamic MPPT efficiency of~\eqref{eq:eta}.

\section{Test profiles used}

The profiles used here follow the medium-to-high class: trapezoidal ramps
between \SI{300}{} and \SI{1000}{\watt\per\square\metre} at \SI{25}{\celsius},
with a \SI{1}{\second} dwell at each plateau and a \SI{1}{\second} settling
period before the first ramp. Four ramp rates were tested: \SI{25}{},
\SI{50}{}, \SI{100}{} and \SI{200}{\watt\per\square\metre\per\second}.

Two deviations from the standard are stated explicitly for honesty. First, one
up-down cycle is used per slope rather than the full repetition count, which
keeps each simulation to a practical length while still exercising both ramp
directions. Second, the fastest rate of \SI{200}{\watt\per\square\metre\per\second}
is above the standard's range, and was included deliberately to push the
trackers past the point where they can keep up. Neither deviation affects the
comparison between trackers, since all four see identical profiles.

\section{Computing the efficiency honestly}

The denominator of~\eqref{eq:eta} requires the true available power at every
instant. Evaluating the exact maximum power point at each of the roughly
\num{500000} simulation samples would be prohibitively slow, so a lookup is
built once: the exact MPP is solved on a grid of \num{240} irradiance values
spanning the profile range, and the result is interpolated with a shape
preserving cubic. Because $\Pmp(G)$ is smooth and nearly linear over this range,
the interpolation error is far below the differences being measured.

The score window excludes the first second, so that the initial acquisition
transient, which is an artefact of the arbitrary starting point, does not
contaminate the tracking measurement.

\section{Results}

\begin{table}[htbp]
\centering
\caption{Dynamic MPPT efficiency (\si{\percent}) on trapezoidal ramps between \SI{300}{} and \SI{1000}{\watt\per\square\metre}.}
\label{tab:dyneta}
\begin{tabular}{@{}lrrrr@{}}
\toprule
\textbf{Tracker} & \SI{25}{} & \SI{50}{} & \SI{100}{} & \SI{200}{\watt\per\square\metre\per\second} \\
\midrule
P\&O, fixed step        & 99.952 & 99.930 & 99.871 & 99.483 \\
P\&O, variable step     & 99.836 & 98.726 & \textbf{93.141} & 99.652 \\
Incremental conductance & 99.950 & 99.926 & 99.871 & 99.458 \\
Extremum seeking        & 99.851 & 99.851 & 99.851 & \textbf{99.850} \\
\bottomrule
\end{tabular}
\end{table}

\begin{table}[htbp]
\centering
\caption{Worst instantaneous power shortfall (\si{\watt}) on the same tests, against a peak available power of \SI{2402}{\watt}.}
\label{tab:dyndef}
\begin{tabular}{@{}lrrrr@{}}
\toprule
\textbf{Tracker} & \SI{25}{} & \SI{50}{} & \SI{100}{} & \SI{200}{\watt\per\square\metre\per\second} \\
\midrule
P\&O, fixed step        & 6.24  & 8.21   & 23.37  & 116.17 \\
P\&O, variable step     & 43.58 & 240.29 & 704.85 & 40.50 \\
Incremental conductance & 6.84  & 9.15   & 17.93  & 129.26 \\
Extremum seeking        & \textbf{6.43} & \textbf{6.63} & \textbf{7.04} & \textbf{7.98} \\
\bottomrule
\end{tabular}
\end{table}

\begin{figure}[htbp]
\centering
\includegraphics[width=\textwidth]{s05_dynamic_benchmark.png}
\caption[Dynamic MPPT benchmark]{Dynamic benchmark results. Top left, efficiency against ramp rate; extremum seeking is flat while the hill-climbers degrade, and the variable-step P\&O collapses at \SI{100}{\watt\per\square\metre\per\second}. Top right, the test profile. Bottom left, voltage tracking, where the runaway excursion of the variable-step tracker down past \SI{90}{\volt} is plainly visible. Bottom right, instantaneous power shortfall.}
\label{fig:benchmark}
\end{figure}

\section{Discussion}

\subsection{Extremum seeking is essentially rate-independent}

The most striking result in Table~\ref{tab:dyneta} is the ESC row: \SI{99.851}{},
\SI{99.851}{}, \SI{99.851}{}, \SI{99.850}{\percent}. An eightfold change in ramp
rate produces a change of \SI{0.001}{\percent}. Its worst instantaneous
shortfall rises only from \SI{6.4}{} to \SI{8.0}{\watt}, whereas the
hill-climbers rise to \SI{116}{} and \SI{129}{\watt}.

The reason is structural. ESC estimates the gradient continuously from a dither
at a known frequency, and the washout filter explicitly rejects the slowly
varying component of power. A slow irradiance ramp is precisely such a slowly
varying component, so ESC is largely blind to it by construction. The
hill-climbers, by contrast, sample power at \SI{50}{\hertz} and cannot
distinguish a power change caused by their own perturbation from one caused by
the ramp.

This is a genuinely control-theoretic advantage: frequency separation between
the probe and the disturbance, rather than a better search heuristic.

\subsection{The cost of that immunity}

ESC pays for this with its dither. Under constant irradiance it scores
\SI{99.55}{\percent} against \SI{99.94}{\percent} for fixed-step P\&O
(Table~\ref{tab:stepcompare}), because the \SI{6}{\volt} peak-to-peak excursion
sits permanently on a curved P--V characteristic and the average power on an arc
is below the power at its centre. The optimal choice therefore depends on the
site: ESC for locations with frequent fast irradiance transients, P\&O for
stable clear-sky conditions.

\subsection{Fixed-step P\&O and INC remain nearly identical}

The two classical hill-climbers agree to within \SI{0.03}{\percent} at every
ramp rate. As argued in Chapter~\ref{ch:mppt} this is expected: with the same
step and sample period they compute the sign of the same quantity by two routes.
INC's advantages are real but appear only in the presence of measurement noise
and quantisation, which this idealised simulation does not include. Adding
sensor noise models is identified as future work in
Chapter~\ref{ch:conclusions}.

\subsection{The variable-step failure}

The variable-step P\&O, which was the \emph{best} tracker under the step test,
is by far the worst here, dropping to \SI{93.14}{\percent} at
\SI{100}{\watt\per\square\metre\per\second} with a peak shortfall of
\SI{705}{\watt}, nearly a third of the available power.

The mechanism is the drift problem of Section~\ref{sec:drift} amplified by the
adaptive rule. The step size is made proportional to $|\Delta P / \Delta V|$,
which is intended to shrink the step near the flat peak and enlarge it on the
steep flanks. During a rising ramp, though, the algorithm locks into the wrong
direction and starts moving away from the peak. As it does so
$|\mathrm{d}P/\mathrm{d}V|$ \emph{grows}, so the adaptive rule enlarges the
step, which moves it further away faster. The feedback is positive. The bottom
left panel of Figure~\ref{fig:benchmark} shows the excursion running from
\SI{158}{\volt} down past \SI{90}{\volt} before the ramp ends and the algorithm
recovers.

The non-monotonic efficiency across ramp rates, worst at \SI{100}{} but
acceptable at \SI{200}{\watt\per\square\metre\per\second}, is consistent with
this explanation: whether a runaway happens depends on the phase relationship
between the perturbation sequence and the ramp, not smoothly on the ramp rate.

This result is reported rather than tuned away because it is instructive. It
demonstrates that an algorithm modification which improves the usual benchmark
can introduce a serious failure mode elsewhere, and that positive feedback
inside an adaptation rule deserves the same stability scrutiny as the main
control loop. The standard remedy in the literature is to add an irradiance
change detector, which tests whether a large $\Delta P$ was accompanied by a
correspondingly small $\Delta V$ and, if so, suppresses the direction reversal.
Implementing and validating such a drift-free variant is the most immediately
valuable extension to this project.

\chapter{Software Architecture and Workflow}
\label{ch:software}

\section{The workflow end to end}

The project was built as a pipeline in which each stage is verified before the
next depends on it. Figure~\ref{fig:workflow} shows the whole chain.

\begin{figure}[htbp]
\centering
\begin{tikzpicture}[
  font=\small,
  node distance=6mm,
  stage/.style={draw, rounded corners=2pt, minimum height=9mm, text width=36mm, align=center},
  check/.style={draw, dashed, rounded corners=2pt, minimum height=8mm, text width=36mm, align=center, font=\footnotesize},
  >={Stealth[length=2.2mm]}
]
\node[stage] (ds) {Manufacturer datasheet\\ \footnotesize $\Voc,\Isc,\Vmp,\Imp,\beta_{Voc}$};
\node[stage, below=of ds] (fit) {Five-parameter identification\\ \footnotesize simultaneous nonlinear solve};
\node[stage, below=of fit] (pv) {PV array model\\ \footnotesize analytic $\mathrm{d}I/\mathrm{d}V$, exact MPP};
\node[stage, below=of pv] (avg) {Averaged converter model\\ \footnotesize state-space, linearised};
\node[stage, below=of avg] (ctrl) {Cascaded control design\\ \footnotesize margins over the envelope};
\node[stage, below=of ctrl] (mppt) {MPPT layer\\ \footnotesize P\&O, INC, ESC};
\node[stage, below=of mppt] (bench) {Dynamic benchmark\\ \footnotesize EN 50530-style ramps};
\node[stage, below=of bench] (sw) {Switched Simscape model\\ \footnotesize real devices at \SI{20}{\kilo\hertz}};
\node[stage, below=of sw] (hw) {Hardware plan\\ \footnotesize ratings, sensing, protection};

\foreach \a/\b in {ds/fit, fit/pv, pv/avg, avg/ctrl, ctrl/mppt, mppt/bench, bench/sw, sw/hw}
  \draw[->] (\a) -- (\b);

\node[check, right=14mm of fit] (c1) {residual at machine precision};
\node[check, right=14mm of pv]  (c2) {slope vs finite difference $<10^{-5}$};
\node[check, right=14mm of avg] (c3) {Jacobian vs finite difference $<10^{-9}$};
\node[check, right=14mm of ctrl](c4) {phase margin $>\SI{45}{\degree}$ everywhere};
\node[check, right=14mm of mppt](c5) {limit cycle $=2\Delta V$};
\node[check, right=14mm of sw]  (c6) {switched vs averaged $<\SI{0.25}{\percent}$};

\foreach \a/\b in {fit/c1, pv/c2, avg/c3, ctrl/c4, mppt/c5, sw/c6}
  \draw[->, dashed] (\a) -- (\b);
\end{tikzpicture}
\caption[Project workflow]{The project workflow. Solid arrows are the modelling chain; dashed boxes are the verification applied at each stage before the next one is allowed to depend on it.}
\label{fig:workflow}
\end{figure}

The organising principle is that \textbf{no stage is trusted until it is checked
against something independent}. The parameter fit is checked against the
datasheet; the analytic derivatives against finite differences; the averaged
model against a switched circuit simulation; the simplified design models
against the full-order plant. This matters because the errors that are dangerous
in a study like this are not crashes but plausible wrong numbers.

\section{Repository layout}

\begin{table}[htbp]
\centering
\caption{Repository structure.}
\label{tab:repo}
\small
\begin{tabular}{@{}>{\raggedright\arraybackslash}p{0.40\textwidth}>{\raggedright\arraybackslash}p{0.52\textwidth}@{}}
\toprule
\textbf{Path} & \textbf{Contents} \\
\midrule
\texttt{matlab/+mpptlab/+pv/} & \texttt{PVArray}, \texttt{extractParams}, \texttt{datasheet} \\
\texttt{matlab/+mpptlab/+conv/} & \texttt{Boost}: averaging, linearisation, envelope \\
\texttt{matlab/+mpptlab/+ctrl/} & \texttt{designCascade}, \texttt{loopAnalysis}, \texttt{PI}, \texttt{designSingleLoopPI} \\
\texttt{matlab/+mpptlab/+mppt/} & \texttt{Tracker}, \texttt{PandO}, \texttt{IncCond}, \texttt{ESC} \\
\texttt{matlab/+mpptlab/+bench/} & \texttt{simulate}, \texttt{irradiance}, \texttt{trackerSet} \\
\texttt{matlab/+mpptlab/+util/} & configuration, figure style, paths \\
\texttt{matlab/scripts/} & \texttt{s01} to \texttt{s05} analysis scripts \\
\texttt{matlab/tests/} & 27 unit tests in three classes \\
\texttt{matlab/run\_all.m} & runs the whole pipeline \\
\texttt{simulink/+pvsrc/pv\_module.ssc} & PV array as a Simscape component \\
\texttt{simulink/build\_mppt\_model.m} & builds the switched model from code \\
\texttt{simulink/run\_switched\_validation.m} & switched against averaged, closed loop \\
\texttt{results/figures/}, \texttt{results/data/} & generated outputs \\
\texttt{docs/report/} & this report \\
\bottomrule
\end{tabular}
\end{table}

\section{Design decisions worth justifying}

\subsection{A package namespace rather than loose functions}

All code lives under the \texttt{mpptlab} package. Calls read as
\texttt{mpptlab.pv.PVArray} and \texttt{mpptlab.ctrl.loopAnalysis}, which makes
provenance obvious at the call site and prevents name collisions with MATLAB
functions such as \texttt{simulate} and \texttt{config}.

\subsection{Handle classes for physical objects}

\texttt{PVArray} and \texttt{Boost} are handle classes. A PV array is a physical
thing with mutable state, its irradiance and temperature, and passing it by
reference into a simulation loop that updates those properties thousands of
times per second is both natural and efficient.

\subsection{One configuration function}

Every script obtains its array, converter and gains from
\texttt{mpptlab.util.config}. There is exactly one place where the module type,
the array topology, the component values and the crossover frequencies are
defined. Changing the array from six modules to eight is a one-line edit that
propagates correctly through every figure and table.

\subsection{A common tracker interface}

All trackers derive from an abstract \texttt{Tracker} class exposing
\texttt{reset} and \texttt{update}. The base class owns reference clamping and
update counting; subclasses implement only \texttt{advance}. This guarantees the
comparison is fair: every tracker sees the same plant, the same inner loops, the
same sample rate mechanics and the same limits, so measured differences are
attributable to the algorithm and nothing else.

\subsection{The Simulink model is generated by a script}

\texttt{build\_mppt\_model.m} constructs the model block by block and wires it
programmatically. A \texttt{.slx} file is a binary archive: it cannot be
meaningfully diffed, reviewed or merged. A builder script can. The model is
therefore reproducible from source, and a change to a component value or a
control gain appears in version control as a readable one-line change.

\section{Numerical methods}

\subsection{Solving the diode equation quickly}

The implicit diode equation must be solved several times per integration step.
Three measures make the closed-loop simulations tractable: warm starting from
the previous solution, which cuts Newton to two or three iterations; clamping
the exponent argument to avoid overflow during transients; and a backtracking
line search that guarantees the residual magnitude never grows.

\subsection{Integration}

The averaged plant is integrated with fixed-step fourth-order Runge--Kutta at
half the switching period, \SI{25}{\micro\second}, with the discrete controllers
updating every \SI{50}{\micro\second}. A step-size convergence test is part of
the test suite: halving the step changes the final voltage by under
\SI{0.5}{\volt}, confirming the result is not an artefact of the integrator.

\subsection{Scoring}

The available power used in the efficiency denominator comes from the exact MPP
solver evaluated on an irradiance grid and interpolated, as described in
Chapter~\ref{ch:benchmark}.

\section{Verification}

\subsection{The test suite}

Twenty-seven tests run in three classes.

\begin{table}[htbp]
\centering
\caption{Test coverage.}
\label{tab:tests}
\begin{tabular}{@{}p{0.30\textwidth}rp{0.48\textwidth}@{}}
\toprule
\textbf{Class} & \textbf{Tests} & \textbf{What is checked} \\
\midrule
\texttt{tPVModel} & 11 & datasheet reproduction, short and open circuit conditions, MPP stationarity, residual on the whole curve, analytic slope against finite difference, monotonicity, irradiance and temperature trends, array scaling, $r_d = \Vmp/\Imp$ at the MPP, zero-irradiance robustness \\
\texttt{tConverterAndControl} & 9 & operating point is an equilibrium, Jacobian against finite differences, power gain vanishes at the MPP and changes sign across it, modal formulas against eigenvalues, sign of the duty-to-voltage gain, cascade margins above \SI{45}{\degree} over the envelope, single-loop margin below \SI{10}{\degree}, continuous conduction mode \\
\texttt{tMPPTAndBench} & 7 & profile shapes and slopes, closed-loop reference holding against the analytic I--V curve, limit cycle equals twice the step, all trackers converge to the MPP, efficiency bounded by one, integrator step convergence \\
\bottomrule
\end{tabular}
\end{table}

Several of these are not routine. The test that the cascade holds above
\SI{45}{\degree} while the single-loop design falls below \SI{10}{\degree}
encodes the central engineering conclusion of the project as an executable
assertion, so a future change that quietly destroys the argument fails a test.
Likewise the test that the duty-to-power gain vanishes at the MPP pins the
theoretical result of Section~\ref{sec:vanishing}.

\subsection{Reproducing everything}

\begin{table}[htbp]
\centering
\caption{Result of the full end-to-end run.}
\label{tab:runall}
\begin{tabular}{@{}lrr@{}}
\toprule
\textbf{Stage} & \textbf{Status} & \textbf{Time (\si{\second})} \\
\midrule
s01 PV characterisation   & OK   & 14.5 \\
s02 plant analysis        & OK   & 5.8 \\
s03 control design        & OK   & 9.0 \\
s04 MPPT comparison       & OK   & 39.9 \\
s05 dynamic benchmark     & OK   & 356.4 \\
s06 switched validation   & OK   & 156.9 \\
s07 model diagram export  & OK   & 6.3 \\
Unit tests                & 27 / 27 passed & 42.0 \\
\midrule
\textbf{Total}            &      & \textbf{630.8} \\
\bottomrule
\end{tabular}
\end{table}

The command \texttt{run\_all} executes the entire pipeline, regenerates every
figure and table in this report, rebuilds the Simulink model from its builder
script, and runs the test suite. Table~\ref{tab:runall} is the result of that
run on the development machine. Every number quoted anywhere in this report
comes from that pipeline; none were transcribed by hand from an earlier
experiment.

\chapter{Physical Model and Hardware}
\label{ch:physical}

\section{Purpose of the physical model}

Everything in Chapters~\ref{ch:converter} to~\ref{ch:benchmark} rests on the
averaged model~\eqref{eq:av1}--\eqref{eq:av2}, which replaces a transistor
switching \num{20000} times a second with a continuous duty ratio. That
replacement is an approximation, and an approximation that has not been checked
is an assumption.

The physical model exists to check it, and to be the bridge to hardware. It
contains a real circuit: a transistor that is either on or off, a diode that
conducts in one direction, an inductor with resistance, capacitors that charge
and discharge every switching period, and a photovoltaic array that is a genuine
two-terminal electrical device rather than a formula.

\section{The photovoltaic array as a Simscape component}

\subsection{Why not a controlled current source}

The obvious way to model a panel in Simulink is a controlled current source
driven by a block that computes $\ipv$ from the sensed $\vpv$. That was the
first approach tried here, and it fails: Simulink reports an algebraic loop.

The reason is instructive. The sensed voltage feeds a block with direct
feedthrough, which drives a current source in the same electrical network the
voltage was sensed from. Simulink cannot see that $\vpv$ is a capacitor state
and therefore not instantaneously dependent on the injected current, so it
conservatively declares a loop. Breaking it requires inserting an artificial
delay or filter, which introduces an error that has nothing to do with the
physics.

\subsection{The component actually used}

The correct solution is to make the array a native Simscape component so the
implicit diode equation becomes part of the differential-algebraic system the
Simscape solver already solves. A component was written in the Simscape language
declaring two electrical nodes, two physical signal inputs for irradiance and
cell temperature, the nine reference parameters, and the constitutive equations:

\begin{verbatim}
branches
    i : n.i -> p.i;
end

equations
    v == p.v - n.v;
    i == Iph - I0*(exp(vd/aa) - 1) - vd/Rsh;
end
\end{verbatim}

with the De Soto translation~\eqref{eq:tr_a}--\eqref{eq:tr_r} computed in a
\texttt{let} block from the inputs. The branch is declared from the negative to
the positive node so that a positive current is one delivered to the external
circuit, matching the sign convention used throughout the report.

The benefits are substantial. There is no algebraic loop, because the solver
handles the implicit equation natively. There is no Newton iteration written by
hand inside a Simulink block. The units are checked at compile time. And the
component is a reusable library block: it appears in a generated library and can
be dropped into any Simscape circuit.

\section{The switched model}

\begin{figure}[htbp]
\centering
\includegraphics[width=\textwidth]{s07_simulink_model.png}
\caption[Simulink switched model]{The switched Simulink and Simscape model. Left, the photovoltaic array component driven by irradiance and cell temperature. Centre, the boost power stage with current and voltage sensing. Bottom, the discrete cascaded controller, the one-sample computation delay, the sawtooth carrier and the pulse width modulation comparator.}
\label{fig:slxmodel}
\end{figure}

\subsection{Power stage}

The circuit of Figure~\ref{fig:boostckt} is built from Simscape Foundation
Library electrical blocks: capacitors, an inductor, a resistor for the winding
loss, an ideal switch as the transistor, a diode, a DC voltage source for the
bus, and current and voltage sensors. Notably, no Simscape Electrical licence is
required; the Foundation Library and one custom component are sufficient.

\subsection{Pulse width modulation}

The duty ratio produced by the controller is compared against a sawtooth carrier
running at \SI{20}{\kilo\hertz} that ramps from \num{0} to \num{1} each period.
The comparison output drives the switch. This is genuine pulse width modulation
with genuine switching instants, not an averaged approximation.

\subsection{The one-sample computation delay}

The controller is a discrete block running at the switching period. Connecting
its output directly to the comparator produced a second algebraic loop, because
the gate state affects the sensed currents instantaneously through the resistive
path while the controller has direct feedthrough from those currents.

The fix adopted is not a numerical trick but a model refinement: a unit delay
between the controller output and the comparator. A real digital controller
samples at the start of a period, computes, and applies the resulting duty ratio
in the \emph{next} period. That one-sample computation delay is a real property
of every digital power converter, and including it makes the model more faithful
rather than less. At \SI{50}{\micro\second} it contributes about
\SI{36}{\degree} of phase lag at the \SI{2}{\kilo\hertz} inner crossover, which
is precisely why the \SI{85}{\degree} phase margin designed in
Chapter~\ref{ch:control} matters: the design has enough headroom to absorb it.

\section{Validation against the averaged model}

\subsection{Open loop at fixed duty}

With the duty ratio held at the analytically predicted \num{0.6112} and no
feedback at all, the switched circuit was allowed to settle from zero initial
conditions, and its waveforms were averaged over the final \SI{20}{\percent} of
the run.

\begin{table}[htbp]
\centering
\caption{Switched circuit against the averaged model at fixed duty \num{0.6112}, \SI{1000}{\watt\per\square\metre}, \SI{25}{\celsius}.}
\label{tab:swvalid}
\begin{tabular}{@{}lrrr@{}}
\toprule
\textbf{Quantity} & \textbf{Switched} & \textbf{Averaged} & \textbf{Difference} \\
\midrule
$\vpv$ mean (\si{\volt})            & 158.169 & 157.800 & \SI{+0.234}{\percent} \\
$i_L$ mean (\si{\ampere})           & 15.198  & 15.220  & \SI{-0.145}{\percent} \\
$i_L$ ripple (\si{\ampere})         & 3.176   & 3.215   & \SI{-1.23}{\percent} \\
$\vpv$ ripple (\si{\volt})          & 0.1946  & 0.2009  & \SI{-3.1}{\percent} \\
$P_{\mathrm{pv}}$ mean (\si{\watt}) & 2401.60 & 2401.72 & \SI{-0.005}{\percent} \\
\bottomrule
\end{tabular}
\end{table}

The agreement in Table~\ref{tab:swvalid} is close, and the residual differences
are all explainable. The mean voltage is \SI{0.23}{\percent} high because the
averaged model neglects the transistor on-resistance and the diode forward drop,
which the switched model includes; those losses slightly reduce the current the
converter draws, so the operating point sits marginally higher on the I--V
curve. The ripple predictions from~\eqref{eq:iripple} are within
\SI{1.3}{\percent}, and the mean power agrees to \SI{0.005}{\percent}.

This validates the entire analytical chain. The linearisation, the frequency
responses, the stability margins and the algorithm comparison all inherit their
credibility from the averaged model, and the averaged model has now been checked
against a circuit that makes none of its assumptions.

\subsection{Closed loop}

With the full cascaded controller and a perturb and observe tracker running on
the switched circuit, through a \SI{1000}{} to \SI{500}{\watt\per\square\metre}
step, the measured MPPT efficiency is \SI{99.93}{\percent}, consistent with the
\SI{99.94}{\percent} obtained on the averaged model in
Table~\ref{tab:stepcompare}.

\begin{figure}[htbp]
\centering
\includegraphics[width=\textwidth]{s06_switched_validation.png}
\caption[Switched circuit validation]{Switched circuit results. Top left and top right, the \SI{20}{\kilo\hertz} ripple in inductor current and PV voltage, with the averaged model prediction as a dashed line. Bottom left, closed-loop tracking with the staircase perturb and observe reference visible. Bottom right, harvested power against the available maximum through the irradiance step.}
\label{fig:swvalid}
\end{figure}

\section{Hardware requirements}

This section specifies what would be needed to build the system that has been
modelled. Ratings are derived from the simulated operating envelope with
engineering margin.

\subsection{Electrical stresses to design against}

\begin{table}[htbp]
\centering
\caption{Worst-case electrical stresses from the simulated envelope.}
\label{tab:stress}
\begin{tabular}{@{}p{0.40\textwidth}rl@{}}
\toprule
\textbf{Quantity} & \textbf{Value} & \textbf{Origin} \\
\midrule
Maximum array open-circuit voltage & \SI{197.4}{\volt} & STC, 6 modules in series \\
Open-circuit voltage at \SI{-10}{\celsius} & $\approx$ \SI{223}{\volt} & cold-morning worst case \\
Maximum array current & \SI{16.4}{\ampere} & $\Isc$ at STC \\
Peak inductor current & \SI{16.8}{\ampere} & mean plus half ripple \\
Switch and diode blocking voltage & \SI{400}{\volt} & DC bus \\
Maximum throughput power & \SI{2402}{\watt} & STC \\
Switching frequency & \SI{20}{\kilo\hertz} & design choice \\
\bottomrule
\end{tabular}
\end{table}

The cold open-circuit voltage matters and is easy to overlook. With
$\beta_{Voc} = \SI{-0.123}{\volt\per\kelvin}$ per module, six modules in series
at \SI{-10}{\celsius} reach roughly \SI{223}{\volt}, some \SI{13}{\percent}
above the STC figure. Input-side components must be rated for that, not for the
STC number.

\subsection{Component selection}

\begin{table}[htbp]
\centering
\caption{Proposed component selection with justification.}
\label{tab:bom}
\begin{tabular}{@{}p{0.20\textwidth}p{0.30\textwidth}p{0.42\textwidth}@{}}
\toprule
\textbf{Item} & \textbf{Specification} & \textbf{Reason} \\
\midrule
Transistor & \SI{650}{\volt} silicon carbide MOSFET, $R_{ds(on)} \le \SI{40}{\milli\ohm}$, \SI{30}{\ampere} & \SI{400}{\volt} bus plus switching overshoot needs \SI{650}{\volt}; silicon carbide keeps switching loss low at \SI{20}{\kilo\hertz} \\
Diode & \SI{650}{\volt} silicon carbide Schottky, \SI{20}{\ampere} & zero reverse recovery removes the dominant turn-on loss and the associated electromagnetic interference \\
Inductor & \SI{1.5}{\milli\henry}, \SI{20}{\ampere} saturation, $\le \SI{0.15}{\ohm}$, powdered iron or ferrite with gap & value set by the \SI{20}{\percent} ripple target; saturation rating above the \SI{16.8}{\ampere} peak \\
Input capacitor & \SI{100}{\micro\farad}, \SI{450}{\volt}, film or low-equivalent-series-resistance electrolytic & sets PV voltage ripple; film preferred for lifetime and ripple current capability \\
Output capacitor & \SI{470}{\micro\farad}, \SI{450}{\volt} & bus stiffening and ripple current \\
Gate driver & isolated, $\ge \SI{4}{\ampere}$ peak, with desaturation protection & the switch source is at circuit ground here, but isolation eases layout and protects the controller \\
Controller & microcontroller or digital signal processor with $\ge \SI{100}{MIPS}$, two channels of \SI{12}{bit} analogue to digital conversion, one pulse width modulation timer & the control law is light; the requirement is deterministic \SI{50}{\micro\second} timing \\
\bottomrule
\end{tabular}
\end{table}

\subsection{Sensing}

The controller needs three measurements.

\begin{itemize}[leftmargin=1.4em]
  \item \term{PV voltage}: a resistive divider scaling \SI{0}{} to
        \SI{250}{\volt} onto the analogue-to-digital converter range, with an
        anti-alias filter. A corner around \SI{5}{\kilo\hertz} is a reasonable
        compromise: high enough not to intrude on the \SI{2}{\kilo\hertz} inner
        loop, low enough to attenuate switching noise.
  \item \term{PV current}: a low-side shunt with a current-sense amplifier, or a
        Hall-effect sensor if galvanic isolation is wanted. Resolution matters
        more than absolute accuracy, because the tracker works on
        \emph{differences}.
  \item \term{Inductor current}: required for the inner loop, and best sampled
        synchronously at the centre of the switch on-time, where the
        instantaneous value equals the cycle average. This standard technique
        removes switching ripple from the measurement without a filter and
        without added phase lag.
\end{itemize}

Sensor resolution deserves emphasis. Near the maximum power point the power
gradient is by definition zero, so the power differences a hill-climbing tracker
must resolve are small. A \SI{1.5}{\volt} perturbation at
\SI{2400}{\watt} changes power by only a few watts, roughly
\SI{0.1}{\percent}. Twelve-bit conversion on both channels is the practical
minimum, and averaging several samples per control period is advisable.

\subsection{Protection}

\begin{itemize}[leftmargin=1.4em]
  \item Input fuse or circuit breaker rated above $\Isc$ with photovoltaic-rated
        breaking capability, since a photovoltaic source cannot supply the large
        fault current an ordinary fuse expects.
  \item Cycle-by-cycle overcurrent trip in the pulse width modulation hardware,
        independent of software.
  \item DC bus overvoltage trip, protecting against the load disappearing.
  \item Overtemperature shutdown on the heatsink.
  \item Anti-windup and duty limiting in software, already implemented in the
        control law, plus a watchdog timer.
\end{itemize}

\subsection{Thermal design}

At full power the transistor conducts roughly \SI{16}{\ampere} for
\SI{61}{\percent} of the period, giving a conduction loss near \SI{6}{\watt} at
\SI{40}{\milli\ohm}, with switching loss of a similar order at
\SI{20}{\kilo\hertz}. The diode conducts \SI{16}{\ampere} for \SI{39}{\percent}
of the period at around \SI{1.5}{\volt}, roughly \SI{9}{\watt}. The inductor
dissipates $I^2 R \approx \SI{35}{\watt}$ at \SI{0.15}{\ohm}, which is the
dominant loss and argues for a lower-resistance winding in a real build. A
heatsink of about \SI{2}{\kelvin\per\watt} with modest airflow is adequate for
the semiconductors.

\section{Staged build plan}

\begin{table}[htbp]
\centering
\caption{Proposed hardware build sequence, each stage gated by a measurable outcome.}
\label{tab:buildplan}
\begin{tabular}{@{}rp{0.32\textwidth}p{0.46\textwidth}@{}}
\toprule
\textbf{Stage} & \textbf{Activity} & \textbf{Exit criterion} \\
\midrule
1 & Controller board bring-up, no power stage & Pulse width modulation at \SI{20}{\kilo\hertz} with commanded duty; analogue channels reading correctly \\
2 & Open-loop power stage on a DC supply with a resistive source emulating the array & Measured voltage conversion ratio matches~\eqref{eq:dutyexact}; ripple matches~\eqref{eq:iripple} within \SI{10}{\percent} \\
3 & Close the inner current loop & Step in current reference settles in under \SI{1}{\milli\second} with no overshoot; measured loop gain crosses near \SI{2}{\kilo\hertz} \\
4 & Close the outer voltage loop & \SI{140}{} to \SI{165}{\volt} reference step settles within \SI{15}{\milli\second}, zero steady-state error \\
5 & Substitute a photovoltaic array simulator for the DC supply & Measured I--V curve matches the identified model; static MPPT efficiency above \SI{99}{\percent} \\
6 & Enable the MPPT layer & Tracker converges from any starting point; limit cycle equals twice the step \\
7 & Run EN 50530 ramp profiles on the simulator & Dynamic MPPT efficiency within \SI{0.5}{\percent} of the simulated values in Table~\ref{tab:dyneta} \\
8 & Outdoor testing with real modules & Sustained operation across a day; energy harvested compared against a reference sweep \\
\bottomrule
\end{tabular}
\end{table}

Two points about this plan. First, a \term{photovoltaic array simulator}, a
programmable supply that reproduces a specified I--V curve, is what makes stages
5 to 7 repeatable; real sunlight is not a controlled input and cannot be used to
compare algorithms. Second, the exit criteria are quantitative and taken
directly from the simulation results in this report, so the hardware campaign
tests the model rather than merely producing waveforms.

\section{Limitations of the physical model}

Stated plainly, so that the results are not over-claimed:

\begin{itemize}[leftmargin=1.4em]
  \item The DC bus is treated as an ideal stiff source. A real bus has finite
        capacitance and its own regulator, which introduces additional dynamics.
  \item Switching devices are ideal apart from on-resistance and forward drop.
        Turn-on and turn-off transitions, parasitic inductance, reverse recovery
        and gate drive delay are not modelled, so this model cannot be used for
        electromagnetic interference or switching loss prediction.
  \item Cell temperature is an input rather than a state. A real panel has
        thermal mass and its temperature follows irradiance with a time constant
        of minutes.
  \item Uniform illumination is assumed, so no partial shading and no multiple
        local maxima.
  \item Sensor noise, quantisation and offset are not modelled, which is why
        incremental conductance shows no advantage over perturb and observe in
        Chapter~\ref{ch:benchmark}.
\end{itemize}

\chapter{Conclusions}
\label{ch:conclusions}

\section{What was required and what was delivered}

The project asked for feedback control to hold a photovoltaic array near maximum
power under changing irradiance, with a MATLAB photovoltaic model, an MPPT
model and a control analysis, and a physical model containing the panel and the
DC--DC converter. Table~\ref{tab:delivered} records what was produced against
each of those.

\begin{table}[htbp]
\centering
\caption{Requirements against deliverables.}
\label{tab:delivered}
\begin{tabular}{@{}p{0.26\textwidth}p{0.50\textwidth}p{0.16\textwidth}@{}}
\toprule
\textbf{Required} & \textbf{Delivered} & \textbf{Evidence} \\
\midrule
Photovoltaic model & Five-parameter single-diode model with De Soto translation; parameters identified from five simultaneous datasheet constraints; analytic derivatives; exact maximum power point solver & Ch.~\ref{ch:pvmodel}, Tables~\ref{tab:pvfit},~\ref{tab:pvvalid} \\
MPPT model & Perturb and observe in fixed and variable step form, incremental conductance and extremum seeking, behind a common discrete-time interface & Ch.~\ref{ch:mppt} \\
Control analysis & Averaged converter model, analytic linearisation verified against finite differences, operating envelope sweeps, cascaded design with margins across the envelope and a single-loop baseline for comparison & Ch.~\ref{ch:converter},~\ref{ch:control} \\
Physical model, panel & Custom Simscape component solving the implicit diode equation natively & Ch.~\ref{ch:physical} \\
Physical model, DC--DC & Switched boost converter with real transistor, diode, inductor and capacitors at \SI{20}{\kilo\hertz}, with digital pulse width modulation and computation delay & Ch.~\ref{ch:physical}, Fig.~\ref{fig:slxmodel} \\
Beyond the brief & EN 50530-style dynamic benchmark, 27 automated tests, scripted reproducible pipeline, hardware specification and staged build plan & Ch.~\ref{ch:benchmark},~\ref{ch:software},~\ref{ch:physical} \\
\bottomrule
\end{tabular}
\end{table}

\section{Principal findings}

\subsection{The plant is an operating-point-dependent resonance}

The dynamic resistance of the array appears directly in the damping term of the
converter's characteristic polynomial, equation~\eqref{eq:zeta}. Along the I--V
curve at fixed irradiance the resonant frequency moves by under
\SI{4}{\percent} while the damping ratio changes by a factor of \num{44}, from
\num{0.023} to \num{1.035}. The array is, in effect, the damping element of its
own converter, and it stops damping precisely when it is behaving as a current
source.

\subsection{Regulating the input-side voltage avoids the right half plane zero}

The transfer function from duty ratio to photovoltaic voltage,
equation~\eqref{eq:Gvd}, has a constant numerator. The non-minimum-phase
behaviour that limits boost converter bandwidth appears only in the
duty-to-output path. Choosing the input-side variable as the controlled output
removes a fundamental bandwidth limitation at no cost.

\subsection{A single proportional--integral loop is structurally inadequate}

Against a resonance peak of \SI{43}{\decibel}, a conventionally designed
single-loop controller produces three gain crossovers and its worst-case phase
margin falls to \SI{1.5}{\degree}. The cascade holds above \SI{85}{\degree}
everywhere, because at the inner crossover the plant reduces to $V_{dc}/(Ls)$
independently of operating point. The improvement is structural, not a matter of
retuning.

\subsection{The loop gain on power vanishes at the optimum}

The linearised gain from duty ratio to photovoltaic power equals
$\mathrm{d}P/\mathrm{d}V$, which is zero at the maximum power point and changes
sign across it. No fixed linear controller can regulate to an extremum. This is
why the MPPT layer must estimate a gradient rather than regulate a value, and it
explains the form of every algorithm in Chapter~\ref{ch:mppt} from first
principles.

\subsection{Timescale separation is a derived design rule}

The measured voltage loop settling time of \SI{11.9}{\milli\second} sets the
minimum MPPT sample period. The \SI{20}{\milli\second} used throughout was
derived from that measurement rather than assumed, and the whole three-layer
hierarchy follows from the same reasoning.

\subsection{Extremum seeking is rate-independent, hill climbing is not}

On trapezoidal ramps, extremum seeking held \SI{99.85}{\percent} at every ramp
rate from \SI{25}{} to \SI{200}{\watt\per\square\metre\per\second}, with a worst
instantaneous shortfall of \SI{8}{\watt}. Perturb and observe and incremental
conductance degraded to \SI{99.46}{\percent} with shortfalls above
\SI{116}{\watt}. The advantage is frequency separation between the probe dither
and the disturbance, and it is paid for by a permanent dither loss that makes
extremum seeking the worst performer under constant irradiance.

\subsection{An adaptive rule can introduce a new instability}

The variable-step perturb and observe tracker was the best performer on the step
test and the worst on the ramp test, collapsing to \SI{93.1}{\percent} with a
\SI{705}{\watt} shortfall. Its step size grows with $|\mathrm{d}P/\mathrm{d}V|$,
so once the drift phenomenon points it the wrong way, moving away from the peak
increases the gradient, which increases the step, which moves it further. The
adaptation loop has positive feedback and needs the same stability scrutiny as
the main control loop.

\subsection{The averaged model is trustworthy}

The switched circuit agrees with the averaged model to \SI{0.23}{\percent} in
mean voltage, \SI{0.15}{\percent} in mean current, \SI{1.2}{\percent} in ripple
and \SI{0.005}{\percent} in mean power. Every analytical conclusion in the report
inherits its credibility from this check.

\section{Limitations}

The limitations of the physical model were listed in Chapter~\ref{ch:physical}.
At the level of the study as a whole, three are worth restating. Uniform
illumination is assumed throughout, so partial shading and the multiple local
maxima it creates are outside scope. No sensor noise, quantisation or offset is
modelled, which flatters both hill-climbing algorithms and removes the main
practical distinction between them. And no hardware was built, so every result
is a simulation result, validated against a more detailed simulation rather than
against measurement.

\section{Future work}

\begin{enumerate}[leftmargin=1.6em]
  \item \textbf{A drift-free perturb and observe variant.} Add an irradiance
        change detector that suppresses the direction reversal when a large
        power change is accompanied by a small voltage change. This directly
        addresses the failure of Section~\ref{sec:drift} and is the single most
        valuable extension.
  \item \textbf{Sensor models.} Add quantisation, noise and offset, which should
        reveal the practical difference between perturb and observe and
        incremental conductance that the idealised study could not.
  \item \textbf{Gain scheduling on dynamic resistance.} The outer loop bandwidth
        varies from \SI{15}{} to \SI{223}{\hertz}. Scheduling the outer
        proportional gain on $r_d$, which the tracker already measures from its
        own perturbations, would equalise it.
  \item \textbf{Ripple correlation control.} Use the converter's own
        \SI{20}{\kilo\hertz} switching ripple as the dither, giving extremum
        seeking behaviour with no added excitation and therefore no dither loss.
        This requires the switched model and is a natural use of the physical
        model already built.
  \item \textbf{Partial shading and global search.} Extend the array model to
        per-module irradiance with bypass diodes, and add a global scan layer
        above the local tracker.
  \item \textbf{Hardware realisation.} Execute the staged plan of
        Table~\ref{tab:buildplan}, using the simulated figures as quantitative
        acceptance criteria.
\end{enumerate}

\section{Closing remark}

The value of framing maximum power point tracking as a control problem is that
the framing is predictive rather than decorative. The vanishing loop gain at the
optimum explains why tracking algorithms must estimate gradients. The dynamic
resistance explains why a naive single loop is fragile and where it is fragile.
The measured settling time sets the tracker sample period. The frequency
separation between dither and disturbance explains why extremum seeking is
immune to ramps. In each case a control-theoretic quantity accounts for
behaviour that would otherwise have to be discovered empirically and accepted
without explanation.


\appendix
\chapter{Parameters and Reference Data}
\label{app:params}

\section{Module datasheet values}

\begin{table}[htbp]
\centering
\caption{Datasheet values at standard test conditions for the two modules in the library.}
\label{tab:datasheets}
\begin{tabular}{@{}lrr@{}}
\toprule
\textbf{Quantity} & \textbf{Kyocera KC200GT} & \textbf{SunPower SPR-305E-WHT-D} \\
\midrule
Technology & multicrystalline silicon & monocrystalline silicon \\
Cells in series $N_s$ & 54 & 96 \\
$\Voc$ (\si{\volt}) & 32.90 & 64.20 \\
$\Isc$ (\si{\ampere}) & 8.21 & 5.96 \\
$\Vmp$ (\si{\volt}) & 26.30 & 54.70 \\
$\Imp$ (\si{\ampere}) & 7.61 & 5.58 \\
$\Pmp$ (\si{\watt}) & 200.14 & 305.23 \\
$\beta_{Voc}$ (\si{\percent\per\kelvin}) & $-0.3739$ & $-0.27269$ \\
$\alpha_{Isc}$ (\si{\percent\per\kelvin}) & $+0.0387$ & $+0.061745$ \\
Fill factor & 0.741 & 0.797 \\
\bottomrule
\end{tabular}
\end{table}

\section{Identified model parameters}

\begin{table}[htbp]
\centering
\caption{Identified reference parameters, module level, at standard test conditions.}
\label{tab:bothfits}
\begin{tabular}{@{}lrr@{}}
\toprule
\textbf{Parameter} & \textbf{KC200GT} & \textbf{SPR-305E} \\
\midrule
$I_{ph,ref}$ (\si{\ampere}) & 8.22713 & 5.96570 \\
$I_{0,ref}$ (\si{\ampere}) & \num{4.3964e-10} & \num{6.4598e-12} \\
$R_{s,ref}$ (\si{\ohm}) & 0.33501 & 0.37350 \\
$R_{sh,ref}$ (\si{\ohm}) & 160.6015 & 393.76 \\
$a_{ref}$ (\si{\volt}) & 1.39246 & 2.33250 \\
Ideality factor $n$ & 1.0036 & 0.9457 \\
Worst datasheet error (\si{\percent}) & \num{1.4e-14} & \num{7.0e-11} \\
\bottomrule
\end{tabular}
\end{table}

\section{Converter and controller parameters}

\begin{table}[htbp]
\centering
\caption{Converter parameters and the design reasoning behind each.}
\label{tab:convparams}
\begin{tabular}{@{}llp{0.46\textwidth}@{}}
\toprule
\textbf{Parameter} & \textbf{Value} & \textbf{Basis} \\
\midrule
Array topology & $6 \times 2$ & \SI{2402}{\watt} at STC, $\Vmp$ below the bus at all irradiances \\
$V_{dc}$ & \SI{400}{\volt} & typical single-phase string inverter bus \\
$f_{sw}$ & \SI{20}{\kilo\hertz} & above audio, low enough for silicon carbide efficiency \\
$L$ & \SI{1.5}{\milli\henry} & \SI{20}{\percent} ripple at the STC maximum power point \\
$C_{\mathrm{pv}}$ & \SI{100}{\micro\farad} & PV voltage ripple below \SI{0.25}{\volt} \\
$C_o$ & \SI{470}{\micro\farad} & bus stiffening \\
$R_L$ & \SI{0.15}{\ohm} & realistic winding resistance for this inductance \\
LC corner & \SI{411}{\hertz} & consequence of $L$ and $C_{\mathrm{pv}}$ \\
Steady duty at STC & \num{0.6112} & equation~\eqref{eq:dutyexact} \\
\bottomrule
\end{tabular}
\end{table}

\begin{table}[htbp]
\centering
\caption{Controller gains.}
\label{tab:gains}
\begin{tabular}{@{}llrl@{}}
\toprule
\textbf{Loop} & \textbf{Gain} & \textbf{Value} & \textbf{Basis} \\
\midrule
Inner current & $K_{p,i}$ & \num{0.04712} & $L\omega_{ci}/V_{dc}$, $f_{ci} = \SI{2}{\kilo\hertz}$ \\
Inner current & $K_{i,i}$ & \num{59.22} & zero a decade below crossover \\
Outer voltage & $K_{p,v}$ & \num{0.15708} & $C_{\mathrm{pv}}\omega_{cv}$, $f_{cv} = \SI{250}{\hertz}$ \\
Outer voltage & $K_{i,v}$ & \num{49.35} & zero a fifth of crossover \\
Single-loop baseline & $K_p$ & \num{2.264e-3} & unity loop gain at \SI{130}{\hertz} \\
Single-loop baseline & $K_i$ & \num{0.3699} & zero at \SI{26}{\hertz} \\
Sample period, both loops & $T_s$ & \SI{50}{\micro\second} & one switching period \\
\bottomrule
\end{tabular}
\end{table}

\section{Tracker parameters}

\begin{table}[htbp]
\centering
\caption{Maximum power point tracker settings used for all comparisons.}
\label{tab:trackparams}
\begin{tabular}{@{}llr@{}}
\toprule
\textbf{Tracker} & \textbf{Parameter} & \textbf{Value} \\
\midrule
All hill climbers & sample period & \SI{20}{\milli\second} \\
Perturb and observe & step & \SI{1.5}{\volt} \\
Perturb and observe, variable & step bounds & \SI{0.25}{} to \SI{6}{\volt} \\
Perturb and observe, variable & adaptation gain & \num{0.06} \\
Incremental conductance & step & \SI{1.5}{\volt} \\
Incremental conductance & tolerance & \num{2e-3} \\
Extremum seeking & sample period & \SI{200}{\micro\second} \\
Extremum seeking & dither frequency & \SI{25}{\hertz} \\
Extremum seeking & dither amplitude & \SI{3}{\volt} \\
Extremum seeking & adaptation gain & \num{20} \\
Extremum seeking & washout and low-pass corner & \SI{5}{\hertz} \\
\bottomrule
\end{tabular}
\end{table}

\section{Physical constants}

\begin{table}[htbp]
\centering
\caption{Physical constants used.}
\label{tab:constants}
\begin{tabular}{@{}llr@{}}
\toprule
\textbf{Constant} & \textbf{Symbol} & \textbf{Value} \\
\midrule
Boltzmann constant & $k$ & \SI{1.380649e-23}{\joule\per\kelvin} \\
Boltzmann constant in electronvolts & $k$ & \SI{8.617333262e-5}{\electronvolt\per\kelvin} \\
Elementary charge & $q$ & \SI{1.602176634e-19}{\coulomb} \\
Silicon bandgap at \SI{25}{\celsius} & $E_{g,ref}$ & \SI{1.121}{\electronvolt} \\
Bandgap temperature coefficient & & \num{-2.677e-4} per kelvin, relative \\
Reference irradiance & $G_{ref}$ & \SI{1000}{\watt\per\square\metre} \\
Reference temperature & $T_{ref}$ & \SI{298.15}{\kelvin} \\
\bottomrule
\end{tabular}
\end{table}

\chapter{Reproducing the Results}
\label{app:repro}

\section{Software required}

MATLAB R2025a with Control System Toolbox, Simulink, Simscape and Optimization
Toolbox. Simscape Electrical is \emph{not} required.

\section{Commands}

\begin{table}[htbp]
\centering
\caption{Commands and what they produce.}
\label{tab:commands}
\begin{tabular}{@{}p{0.36\textwidth}p{0.56\textwidth}@{}}
\toprule
\textbf{Command} & \textbf{Produces} \\
\midrule
\texttt{startup\_mpptlab} & puts the library on the MATLAB path \\
\texttt{run\_all} & the entire pipeline, every figure and table, and the test suite \\
\texttt{s01\_pv\_characterisation} & Figure~\ref{fig:pvchar}, Tables~\ref{tab:pvfit},~\ref{tab:pvvalid},~\ref{tab:pvops} \\
\texttt{s02\_plant\_analysis} & Figure~\ref{fig:plantbode}, Tables~\ref{tab:envG},~\ref{tab:envV} \\
\texttt{s03\_control\_design} & Figure~\ref{fig:controldesign}, Tables~\ref{tab:singleloop},~\ref{tab:cascadeG} \\
\texttt{s04\_mppt\_comparison} & Figure~\ref{fig:mpptcompare}, Tables~\ref{tab:stepcompare},~\ref{tab:tradeoff} \\
\texttt{s05\_dynamic\_benchmark} & Figure~\ref{fig:benchmark}, Tables~\ref{tab:dyneta},~\ref{tab:dyndef} \\
\texttt{run\_switched\_validation} & Figure~\ref{fig:swvalid}, Table~\ref{tab:swvalid} \\
\texttt{export\_model\_diagram} & Figure~\ref{fig:slxmodel} \\
\texttt{runtests('tests')} & the 27 unit tests \\
\bottomrule
\end{tabular}
\end{table}

\section{Runtime}

The full pipeline takes about \SI{11}{\minute}. The dynamic benchmark of
Chapter~\ref{ch:benchmark} accounts for roughly \SI{6}{\minute} of that, because
it simulates \SI{121}{\second} of closed-loop operation for each of four
trackers, and the switched validation for a further \SI{2.6}{\minute}.


\begin{thebibliography}{99}
\addcontentsline{toc}{chapter}{Bibliography}

\bibitem{desoto} W.~De~Soto, S.~A.~Klein and W.~A.~Beckman, ``Improvement and validation of a model for photovoltaic array performance,'' \emph{Solar Energy}, vol.~80, no.~1, pp.~78--88, 2006.

\bibitem{villalva} M.~G.~Villalva, J.~R.~Gazoli and E.~R.~Filho, ``Comprehensive approach to modeling and simulation of photovoltaic arrays,'' \emph{IEEE Transactions on Power Electronics}, vol.~24, no.~5, pp.~1198--1208, 2009.

\bibitem{erickson} R.~W.~Erickson and D.~Maksimovic, \emph{Fundamentals of Power Electronics}, 3rd ed. Springer, 2020.

\bibitem{middlebrook} R.~D.~Middlebrook and S.~Cuk, ``A general unified approach to modelling switching-converter power stages,'' in \emph{IEEE Power Electronics Specialists Conference}, 1976, pp.~18--34.

\bibitem{femia} N.~Femia, G.~Petrone, G.~Spagnuolo and M.~Vitelli, ``Optimization of perturb and observe maximum power point tracking method,'' \emph{IEEE Transactions on Power Electronics}, vol.~20, no.~4, pp.~963--973, 2005.

\bibitem{hussein} K.~H.~Hussein, I.~Muta, T.~Hoshino and M.~Osakada, ``Maximum photovoltaic power tracking: an algorithm for rapidly changing atmospheric conditions,'' \emph{IEE Proceedings -- Generation, Transmission and Distribution}, vol.~142, no.~1, pp.~59--64, 1995.

\bibitem{krstic} M.~Krsti\'{c} and H.-H.~Wang, ``Stability of extremum seeking feedback for general nonlinear dynamic systems,'' \emph{Automatica}, vol.~36, no.~4, pp.~595--601, 2000.

\bibitem{leyva} R.~Leyva, C.~Alonso, I.~Queinnec, A.~Cid-Pastor, D.~Lagrange and L.~Martinez-Salamero, ``MPPT of photovoltaic systems using extremum seeking control,'' \emph{IEEE Transactions on Aerospace and Electronic Systems}, vol.~42, no.~1, pp.~249--258, 2006.

\bibitem{esram} T.~Esram and P.~L.~Chapman, ``Comparison of photovoltaic array maximum power point tracking techniques,'' \emph{IEEE Transactions on Energy Conversion}, vol.~22, no.~2, pp.~439--449, 2007.

\bibitem{en50530} CENELEC, \emph{EN 50530: Overall Efficiency of Grid Connected Photovoltaic Inverters}, 2010.

\bibitem{kyocera} Kyocera Corporation, \emph{KC200GT High Efficiency Multicrystal Photovoltaic Module Datasheet}.

\bibitem{astrom} K.~J.~{\AA}str\"{o}m and R.~M.~Murray, \emph{Feedback Systems: An Introduction for Scientists and Engineers}, 2nd ed. Princeton University Press, 2021.

\end{thebibliography}

\end{document}
